\documentclass[a4paper,11pt]{article}
\usepackage[utf8]{inputenc}
\usepackage[T1]{fontenc}
\usepackage{textcomp}
\usepackage[italian]{babel}
\usepackage{amsmath, amssymb, amsthm}
\usepackage{geometry}
\usepackage{hyperref}
\usepackage{xcolor}
\usepackage{tcolorbox}
\usepackage{enumitem}
\usepackage{booktabs}
\usepackage{algorithm}
\usepackage{algpseudocode}
\usepackage[protrusion=true,expansion=false]{microtype}
\usepackage{tikz}
\usepackage{capt-of}
\usepackage{pgfplots}
\pgfplotsset{compat=1.18}
\usepgfplotslibrary{fillbetween}
\usetikzlibrary{arrows.meta, positioning, shapes.geometric, shapes.multipart, calc, fit, backgrounds, decorations.pathreplacing, decorations.markings, trees, angles, quotes}

\geometry{margin=2.5cm}

\tcbuselibrary{theorems, skins, breakable}

\newtcolorbox{definizione}[1][]{colback=red!5!white, colframe=red!60!black,
  fonttitle=\bfseries, title={Definizione: #1}, breakable}
\newtcolorbox{teorema}[1][]{colback=yellow!5!white, colframe=orange!80!black,
  fonttitle=\bfseries, title={Legge/Teorema: #1}, breakable}
\newtcolorbox{esempio}[1][]{colback=green!5!white, colframe=green!50!black,
  fonttitle=\bfseries, title={Esempio: #1}, breakable}
\newtcolorbox{nota}[1][]{colback=blue!5!white, colframe=blue!40!black,
  fonttitle=\bfseries, title={Nota: #1}, breakable}

\hypersetup{colorlinks=true, linkcolor=blue, urlcolor=blue}

\title{\textbf{Fisica per Applicazioni di Realtà Virtuale}\\[0.5em]
\large Note Complete del Corso\\[0.3em]
\normalsize Università di Torino -- Magistrale Informatica -- Prof. Matteo Brogi}
\author{}
\date{}

\begin{document}
\maketitle
\tableofcontents
\newpage

%=============================================================================
\section{Introduzione e Concetti Fondamentali}
%=============================================================================

\subsection{La Fisica come Scienza Sperimentale}

\begin{definizione}[Fisica]
La \textbf{fisica} \`{e} lo studio dei fenomeni naturali; il \emph{come} si studiano tali fenomeni \`{e} la \textbf{misura delle grandezze fisiche}. Alla base della fisica c'\`{e} sempre l'\textbf{evidenza sperimentale}.
\end{definizione}

\begin{center}
\begin{tikzpicture}[node distance=1.7cm, font=\small,
  box/.style={draw=teal!70!black, rounded corners, minimum width=2.6cm, minimum height=0.9cm, align=center, fill=teal!5}]
\node[box] (mis) {Misura};
\node[box, right=of mis] (dat) {Dati\\sperimentali};
\node[box, right=of dat] (ana) {Analisi dati\\(statistica)};
\node[box, right=of ana] (mod) {Modello\\interpretativo};
\node[box, below=1.3cm of ana] (ver) {Verifica\\sperimentale};
\draw[-Stealth, thick, red!70!black] (mis) -- (dat);
\draw[-Stealth, thick, red!70!black] (dat) -- (ana);
\draw[-Stealth, thick, red!70!black] (ana) -- (mod);
\draw[-Stealth, thick, red!70!black] (mod) -- (ver);
\draw[-Stealth, thick, red!70!black] (ver) -- (mis);
\end{tikzpicture}
\captionof{figure}{Il ciclo del metodo scientifico: un modello \`{e} \textbf{inadeguato} se non \`{e} verificato dall'evidenza sperimentale; i limiti nelle misure influenzano a loro volta l'evidenza sperimentale disponibile.}
\end{center}

\begin{esempio}[Eratostene e la misura del raggio terrestre (240--230 a.C.)]
Uno dei primi esempi storici di applicazione del metodo scientifico (misura + geometria). Eratostene osserva che, il giorno del solstizio d'estate:
\begin{itemize}
  \item a \textbf{Syene}, il Sole \`{e} perfettamente perpendicolare al suolo (un pozzo profondo \`{e} illuminato fino in fondo, nessun'ombra);
  \item ad \textbf{Alessandria}, distante $d = 5000$ stadi da Syene, i raggi solari formano un angolo di $7.2^{\circ}$ rispetto alla verticale.
\end{itemize}

\begin{center}
\begin{tikzpicture}[font=\small]
\draw[thick] (0,0) circle (2cm);
\coordinate (O) at (0,0);
\coordinate (S) at (0,2);
\coordinate (A) at (2*0.9925,2*0.1219);
\draw[thick, dashed] (O) -- (S);
\draw[thick, dashed] (O) -- (A);
\draw[thick] (O) -- ($(O)!1.6cm!(S)$) node[midway, left]{$R$};
\foreach \x in {-0.6,-0.3,0,0.3,0.6} {
  \draw[-Stealth, orange!80!black] (\x,3.3) -- (\x,3.3-1.3);
}
\node[above] at (0,3.35) {raggi solari (paralleli)};
\draw pic["$7.2^{\circ}$", draw=teal!70!black, angle radius=1.1cm, angle eccentricity=1.4] {angle=A--O--S};
\node[right] at (A) {Alessandria (A)};
\node[above] at (S) {Syene (S)};
\node[below] at (O) {centro Terra};
\end{tikzpicture}
\captionof{figure}{I raggi solari, praticamente paralleli data la distanza Terra--Sole, colpiscono Syene perpendicolarmente e Alessandria con un angolo di $7.2^{\circ}$: per geometria, questo \`{e} anche l'angolo al centro della Terra sotteso dall'arco Alessandria--Syene.}
\end{center}

Per proporzionalit\`{a} tra angolo e arco corrispondente sulla circonferenza:
\[
\frac{\text{circonferenza}}{d} = \frac{360^{\circ}}{7.2^{\circ}} \quad\Rightarrow\quad 2\pi R = d \cdot \frac{360^{\circ}}{7.2^{\circ}} \approx 250\,000 \text{ stadi} \approx 40\,000 \text{ km}
\]
Il valore moderno della circonferenza terrestre \`{e} $40\,100$ km: una precisione notevole per l'epoca, ottenuta con sola misura e geometria -- un esempio ante-litteram del ciclo misura$\to$modello$\to$verifica.
\end{esempio}

\begin{nota}[Due rivoluzioni scientifiche: perch\'{e} questo corso tratta la fisica ``classica'']
\begin{itemize}
  \item \textbf{Prima rivoluzione scientifica} (XVI--XVII sec., Copernico, Galileo, Newton, Cartesio, Keplero\ldots): supera il geocentrismo e i quattro elementi aristotelici, introduce le leggi della dinamica e il metodo scientifico moderno. La \textbf{fisica classica} che ne risulta \emph{predice} il comportamento anche di sistemi complessi, ed \`{e} coerente con l'\textbf{intuizione} dei fenomeni naturali alle scale abituali.
  \item \textbf{Seconda rivoluzione scientifica} (XIX--XX sec., Planck, Einstein, Bohr, Hubble\ldots): meccanica quantistica, relativit\`{a}, fisica atomica -- necessaria per spiegare fenomeni a scale molto piccole o molto grandi, ma \emph{contrasta} con l'intuizione comune; inoltre la misura pu\`{o} essere distruttiva (non ripetibile) e alcuni sistemi vanno descritti solo in modo stocastico.
\end{itemize}
Questo corso si concentra sulla \textbf{fisica classica}: descrive bene la realt\`{a} alle scale abituali, \`{e} coerente con l'intuizione, ed \`{e} quindi il naturale \emph{soggetto} per riprodurre la realt\`{a} virtualmente.
\end{nota}

\subsection{Grandezze Fisiche}

\begin{definizione}[Classificazione delle Grandezze Fisiche]
Le grandezze fisiche si classificano secondo tre coppie dicotomiche indipendenti:
\begin{itemize}
  \item \textbf{Dimensionali} (necessitano di un'unit\`{a} di misura, es. lunghezza, tempo, massa) vs. \textbf{Adimensionali} (numeri puri, es. il \emph{radiante} = arco/raggio).
  \item \textbf{Fondamentali} (massa, spostamento, tempo\ldots) vs. \textbf{Derivate} (combinazione di grandezze fondamentali, es. velocit\`{a}, forza, energia).
  \item \textbf{Scalari} (descritte interamente da un valore) vs. \textbf{Vettoriali} (descritte da un valore/modulo, una direzione e un verso).
\end{itemize}
\end{definizione}

\subsection{Unità di Misura e Analisi Dimensionale}

\begin{definizione}[Sistema Internazionale (SI)]
In uso dal \textbf{1978} per evitare confusione ed errori dovuti a sistemi di unit\`{a} diversi. \textbf{Esempio monito}: la sonda \emph{Mars Climate Orbiter} (1999) and\`{o} persa a causa di un errore di conversione tra unit\`{a} imperiali e metriche in una parte del software di navigazione.

\begin{center}
\small
\begin{tabular}{lll}
\toprule
\textbf{Grandezza fisica} & \textbf{Simbolo} & \textbf{Unit\`{a} (simbolo)} \\
\midrule
Lunghezza & $l$ & metro (m) \\
Massa & $m$ & kilogrammo (kg) \\
Tempo & $t$ & secondo (s) \\
Corrente elettrica & $I$ & ampere (A) \\
Temperatura & $T$ & kelvin (K) \\
Quantit\`{a} di sostanza & $n$ & mole (mol) \\
Intensit\`{a} luminosa & $i_v$ & candela (cd) \\
\bottomrule
\end{tabular}
\end{center}
Le 7 grandezze fondamentali del SI; tutte le altre unit\`{a} (velocit\`{a}, forza, energia\ldots) sono derivate da queste per combinazione.
\end{definizione}

\begin{nota}[Prefissi delle potenze del 10]
Le grandezze fisiche variano su ordini di grandezza molto diversi, da cui la necessit\`{a} di prefissi:
\begin{center}
\small
\begin{tabular}{lccccccc}
\toprule
Potenza & $10^{15}$ & $10^{12}$ & $10^9$ & $10^6$ & $10^3$ & $10^2$ & $10^1$ \\
Prefisso & peta & tera & giga & mega & kilo & etto & deca \\
Simbolo & P & T & G & M & k & h & da \\
\midrule
Potenza & $10^{-1}$ & $10^{-2}$ & $10^{-3}$ & $10^{-6}$ & $10^{-9}$ & $10^{-12}$ & $10^{-15}$ \\
Prefisso & deci & centi & milli & micro & nano & pico & femto \\
Simbolo & d & c & m & \textmu & n & p & f \\
\bottomrule
\end{tabular}
\end{center}
\textbf{Attenzione}: kilo, mega, giga\ldots dovrebbero essere familiari dall'informatica, ma $10^3 \neq 2^{10}$ -- i prefissi SI sono potenze di 10, mentre in informatica si usano spesso potenze di 2 (kibi, mebi, gibi\ldots).

Unit\`{a} di lunghezza non-SI usate in fisica: \r{A}ngstrom ($1\,\text{\r{A}} = 10^{-10}$ m, scala atomica) e parsec ($1\,\text{pc} = 3.09\times10^{16}$ m, scala astronomica).
\end{nota}

\begin{definizione}[Analisi Dimensionale]
Specificare la \textbf{dimensione} di una grandezza significa specificarne il tipo, indicata tra parentesi quadre $[\cdot]$. \`{E} uno strumento essenziale per verificare che una formula abbia senso fisico.
\begin{center}
\small
\begin{tabular}{lc}
\toprule
Grandezza & Dimensione \\
\midrule
Distanza & $[L]$ \\
Area & $[L^2]$ \\
Volume & $[L^3]$ \\
Tempo & $[T]$ \\
Velocit\`{a} & $[L\,T^{-1}]$ \\
Massa & $[M]$ \\
\bottomrule
\end{tabular}
\end{center}
\end{definizione}

\begin{esempio}[Verifica dimensionale]
Verificare la consistenza dimensionale di $x = x_0 + vt$:
\[
[L] = [L] + [L\,T^{-1}][T] = [L] + [L] \quad\checkmark
\]
Ogni addendo ha la stessa dimensione $[L]$: la formula \`{e} dimensionalmente consistente (condizione necessaria, non sufficiente, per la sua correttezza).
\end{esempio}

\subsection{Cifre Significative e Notazione Scientifica}

\begin{definizione}[Cifre Significative]
Ogni misura ha un errore associato, quindi il numero di cifre riportate ha un significato preciso:
\begin{itemize}
  \item Le \textbf{cifre significative} sono il numero di cifre che compongono il numero; le \textbf{cifre decimali} sono il numero di cifre dopo la virgola.
  \item Gli \textbf{zeri all'inizio} del numero non si contano come significativi.
  \item Tutte le cifre sono \textbf{certe} tranne l'\textbf{ultima}, che \`{e} \textbf{incerta}.
\end{itemize}
\end{definizione}

\begin{esempio}[Cifre significative -- $0.01370$]
Gli zeri iniziali ($0.0$) non contano; le cifre certe sono $1,3,7$; la prima cifra incerta \`{e} l'ultimo $0$. Risultato: \textbf{4 cifre significative, 5 cifre decimali}.

\textbf{Caso ambiguo}: il numero $4700$ non specifica se le cifre significative siano 2, 3 o 4 -- l'ambiguit\`{a} si risolve scrivendo il numero in \textbf{notazione scientifica} (sezione seguente).
\end{esempio}

\begin{teorema}[Cifre Significative nelle Operazioni Algebriche]
Regole derivate dalla teoria degli errori:
\begin{itemize}
  \item \textbf{Moltiplicazione/divisione}: il risultato ha tante \textbf{cifre significative} quante la quantit\`{a} pi\`{u} incerta (cio\`{e} quella con meno cifre significative).
  \item \textbf{Somma/sottrazione}: il risultato ha tante \textbf{cifre decimali} quante il termine con il minor numero di cifre decimali.
  \item Si \textbf{arrotonda}, non si tronca (es. $2.55146\ldots \to 2.6$, non $2.5$).
\end{itemize}
\end{teorema}

\begin{esempio}[Cifre significative in pratica]
\textbf{Divisione}: $d = 21.26$ cm (4 cifre sign.), $t = 8.5$ s (2 cifre sign.) $\Rightarrow v = d/t = 2.50117\ldots = \mathbf{2.5\ cm\,s^{-1}}$ (eredita le 2 cifre significative di $t$).

\textbf{Somma}: $v_0 = 1.384$ cm/s (3 decimali), $v = 2.5$ cm/s (1 decimale) $\Rightarrow v+v_0 = \mathbf{3.9\ cm\,s^{-1}}$ (eredita 1 decimale da $v$).
\end{esempio}

\begin{nota}[Errore di arrotondamento nei calcoli multi-passo]
Arrotondare a ogni passaggio intermedio, invece che solo al risultato finale, pu\`{o} introdurre un errore cumulativo.

\textbf{Esempio}: due importi di 2.21\texteuro{} e 1.35\texteuro{}, con tassa dell'8\%.
\begin{itemize}
  \item Arrotondando a ogni passo: $2.21\texteuro \to 2.3868\texteuro \to 2.39\texteuro$; $1.35\texteuro \to 1.458\texteuro \to 1.46\texteuro$; totale $= 2.39+1.46 = \mathbf{3.85\texteuro}$ (impreciso).
  \item Sommando prima, poi applicando la tassa: $(2.21+1.35)\texteuro \times 1.08 = 3.8448\texteuro = \mathbf{3.84\texteuro}$ (corretto).
\end{itemize}
\textbf{Regola pratica}: usare almeno una cifra significativa in pi\`{u} nei calcoli intermedi e arrotondare solo alla fine.
\end{nota}

\begin{definizione}[Notazione Scientifica]
Un numero si esprime con una cifra intera non nulla, seguita da eventuali decimali, moltiplicata per una potenza di 10: $M_\oplus = 5{,}970{,}000{,}000{,}000{,}000{,}000{,}000{,}000$ kg $= 5.97\times10^{24}$ kg (o $5.97E24$).

Risolve l'ambiguit\`{a} di casi come $4700$, perch\'{e} rende esplicito il numero di cifre significative con cui il numero \`{e} scritto (es. $4.7\times10^3$ ha 2 cifre sign., $4.700\times10^3$ ne ha 4).
\end{definizione}

\begin{esempio}[Conversione di unit\`{a}]
Un cartello stradale in miglia (mi) e chilometri (km): 26 mi $\to$ ? km. Fattore di conversione $1\ \text{mi} = 1.609$ km, usato come frazione unitaria:
\[
26\ \text{mi} \times \frac{1.609\ \text{km}}{1\ \text{mi}} = 41.834\ \text{km} \approx 42\ \text{km}
\]
\textbf{Attenzione}: per le \textbf{unit\`{a} derivate} (es. velocit\`{a} = spazio/tempo), ogni unit\`{a} va convertita \textbf{separatamente} (numeratore e denominatore).
\end{esempio}

\begin{nota}[Approssimazione per ordini di grandezza]
Molto comune in fisica (specialmente astrofisica): si approssima alla potenza di 10 pi\`{u} vicina (es. $2\to1$, $8\to10$). Alcune costanti si prestano bene: $\pi^2 \approx 10$; $1$ anno $\approx \pi \times 10^7$ s. Sempre accettabile a meno che l'esercizio non richieda precisione esplicita -- utile per stime rapide (``stime alla Fermi'').
\end{nota}

\subsection{Algebra Vettoriale}

\begin{definizione}[Vettore]
Una quantit\`{a} vettoriale ha un \textbf{modulo}, una \textbf{direzione} e un \textbf{verso}. \textbf{Il punto di applicazione non \`{e} una propriet\`{a} del vettore}: due vettori con lo stesso modulo, direzione e verso sono lo stesso vettore, indipendentemente da dove sono disegnati.
\end{definizione}

\begin{teorema}[Somma e Differenza di Vettori -- Regola del Parallelogramma]
\[
\vec{u}+\vec{v} = \vec{v}+\vec{u} \qquad \vec{u}-\vec{v} = -(\vec{v}-\vec{u})
\]
Somma e differenza sono operazioni \textbf{interne}: il risultato \`{e} un altro vettore.
\end{teorema}

\begin{center}
\begin{tikzpicture}[scale=1.1, font=\small]
\coordinate (A) at (0,0);
\coordinate (B) at (2.5,0.6);
\coordinate (D) at (0.9,2.1);
\coordinate (E) at ($(B)+(D)$);
\draw[-Stealth, thick, blue] (A) -- (B) node[midway, below]{$\vec{u}$};
\draw[-Stealth, thick, teal!70!black] (A) -- (D) node[midway, left]{$\vec{v}$};
\draw[dashed] (B) -- (E) -- (D);
\draw[-Stealth, thick, red!70!black] (A) -- (E) node[midway, above right]{$\vec{u}+\vec{v}$};
\draw[-Stealth, thick, orange!90!black] (D) -- (B) node[midway, above]{$\vec{u}-\vec{v}$};
\end{tikzpicture}
\captionof{figure}{Regola del parallelogramma: la somma $\vec u + \vec v$ \`{e} la diagonale principale, la differenza $\vec u - \vec v$ l'altra diagonale.}
\end{center}

\begin{definizione}[Versori e Componenti Cartesiane]
I \textbf{versori} sono vettori di modulo unitario (solo direzione e verso). I versori cartesiani $\hat{i}, \hat{j}$ (e $\hat{k}$ in 3D) sono perpendicolari tra loro e individuano gli assi. Qualsiasi vettore \`{e} scomponibile come somma di versori cartesiani:
\[
\vec{a} = a_x\hat{i} + a_y\hat{j} \equiv (a_x, a_y)
\]
\end{definizione}

\begin{nota}[In parole semplici: somma vettoriale in componenti]
Grazie alla scomposizione in versori, sommare o sottrarre vettori si riduce a sommare o sottrarre le rispettive componenti:
\[
\vec{a} \pm \vec{b} = (a_x \pm b_x)\,\hat{i} + (a_y \pm b_y)\,\hat{j}
\]
Questo \`{e} esattamente il motivo per cui, in un motore grafico o di fisica (es. per una simulazione VR), i vettori si rappresentano come semplici array di numeri (componenti): sommare due vettori \`{e} solo una somma componente-per-componente, un'operazione immediata da programmare.
\end{nota}

\begin{teorema}[Prodotto Scalare (Dot Product)]
Operazione \textbf{esterna}: il risultato \`{e} uno \textbf{scalare}.
\[
\vec{a}\cdot\vec{b} = ab\cos\theta = a\,b_a \qquad \text{(in componenti: } \vec{a}\cdot\vec{b} = a_xb_x + a_yb_y\text{)}
\]
dove $b_a$ \`{e} la proiezione di $\vec b$ su $\vec a$. \`{E} commutativo: $\vec a \cdot \vec b = \vec b \cdot \vec a$.
\textbf{Casi particolari}: $\vec a \parallel \vec b \Rightarrow \vec a \cdot \vec b = ab$; \quad $\vec a \perp \vec b \Rightarrow \vec a \cdot \vec b = 0$.
\end{teorema}

\begin{esempio}[Il dot product nel rendering 3D -- illuminazione diretta (\emph{direct illumination})]
In computer grafica (rilevante per il rendering di scene in un visore VR), il \textbf{segno} del prodotto scalare tra il versore che punta verso la \textbf{sorgente luminosa} e la \textbf{normale} alla superficie di un triangolo della mesh determina se quel punto \`{e} illuminato o in ombra.

\begin{center}
\begin{tikzpicture}[font=\small]
\draw[thick] (0,0) circle (1.4cm);
\draw[-Stealth, orange, thick] (-2.6,1.8) -- (-1,0.7) node[midway, above]{\footnotesize luce};
\foreach \ang in {40,80,120,160,200,240,290,330} {
  \draw[-Stealth, red!70!black] (\ang:1.4) -- (\ang:2.1);
}
\node[align=center, below] at (-0.6,0.9) {\footnotesize lato illuminato};
\node[align=center, below] at (1.0,-1.0) {\footnotesize lato in ombra};
\end{tikzpicture}
\captionof{figure}{Vettori normali (rosso) uscenti dalla superficie di una mesh sferica; dove la normale punta verso la sorgente di luce (prodotto scalare negativo con la direzione di propagazione della luce) il punto \`{e} illuminato, altrimenti \`{e} in ombra.}
\end{center}
\end{esempio}

\begin{teorema}[Prodotto Vettoriale (Cross Product)]
Operazione ``interna'': il vettore risultante \`{e} \textbf{perpendicolare} a entrambi i vettori di partenza (regola della mano destra per il verso).
\[
|\vec a \times \vec b| = ab\sin\theta \qquad \vec a \times \vec b = -\vec b \times \vec a \quad \text{(anticommutativo)}
\]
\textbf{Casi particolari}: $\vec a \perp \vec b \Rightarrow |\vec a \times \vec b| = ab$; \quad $\vec a \parallel \vec b \Rightarrow \vec a \times \vec b = \vec 0$.

In componenti cartesiane (determinante):
\[
\vec a \times \vec b = \begin{vmatrix} \hat i & \hat j & \hat k \\ a_x & a_y & a_z \\ b_x & b_y & b_z \end{vmatrix}
\]
\end{teorema}

\subsection{Fisica Classica e Simulazione: i Limiti della Discretizzazione}

\begin{nota}[Dalla realt\`{a} continua alla simulazione discreta]
La realt\`{a} fisica (classica) \`{e} \textbf{continua}, mentre qualunque simulazione al calcolatore \`{e} necessariamente \textbf{discreta}:
\begin{itemize}
  \item il tempo \`{e} campionato a intervalli \textbf{finiti} -- ad esempio un frame di un visore VR a \textbf{75 Hz} corrisponde a un intervallo temporale di $1/75\,\text{s} \approx 13.3$ ms;
  \item anche la precisione numerica con cui si rappresentano distanze, velocit\`{a}, ecc. \`{e} \textbf{finita} (aritmetica in virgola mobile).
\end{itemize}
Questo apre la domanda che guida il resto del corso: cosa serve per scrivere codice \emph{efficiente} che simuli correttamente un fenomeno fisico continuo con risorse di calcolo discrete e limitate?
\end{nota}

%=============================================================================
\section{Cinematica}
%=============================================================================

\subsection{Cinematica Unidimensionale}

\begin{definizione}[Cinematica]
La \textbf{cinematica} descrive il moto \emph{senza} trattarne le cause (``come si muove?'', a differenza della dinamica che risponde a ``perch\'{e} si muove?''). Si parte dal moto \textbf{unidimensionale} (una sola direzione, grandezze scalari), \textbf{traslazionale} (niente rotazione), trattando i corpi come \textbf{punti materiali} (oggetti ideali senza dimensione).
\end{definizione}

Su un asse coordinato (orientato, con origine $O$), la posizione di un punto materiale al tempo $t$ è $x$; $t_0$ è il tempo di riferimento.

\begin{definizione}[Velocit\`{a} Scalare]
\[
v_m = \bar v = \frac{\Delta x}{\Delta t} = \frac{x_2-x_1}{t_2-t_1} \quad \text{(media)} \qquad\qquad v = \lim_{\Delta t \to 0}\frac{\Delta x}{\Delta t} = \frac{dx}{dt} \quad \text{(istantanea)}
\]
La velocit\`{a} istantanea ha un \textbf{significato geometrico}: \`{e} la pendenza della curva $x(t)$ nel punto $(t_1,x_1)$ (retta tangente).
\end{definizione}

\begin{center}
\begin{tikzpicture}
\begin{axis}[width=8.5cm, height=6cm, xlabel={$t$}, ylabel={$x$}, xtick=\empty, ytick=\empty, axis lines=left, font=\small]
\addplot[domain=0:4, samples=100, thick, blue] {0.3*x^2+0.2};
\addplot[mark=*, mark size=2pt, color=red] coordinates {(2,1.4)};
\addplot[domain=0.8:3.2, thick, dashed, green!50!black] {1.2*x-1};
\node[anchor=west] at (axis cs:2,1.4) {\ $(t_1,x_1)$};
\node[anchor=west, green!50!black] at (axis cs:1.5,0.9) {\footnotesize pendenza $=v_1$};
\end{axis}
\end{tikzpicture}
\captionof{figure}{La velocit\`{a} istantanea $v_1$ nel punto $(t_1,x_1)$ \`{e} la pendenza della retta tangente al grafico $x(t)$.}
\end{center}

\begin{nota}[Approssimazione discreta della velocit\`{a} istantanea]
In applicazioni software $\Delta t$ \`{e} sempre \textbf{finito}: la scelta del passo temporale deve essere adatta al problema, anche \textbf{adattiva} a seconda della scala del fenomeno (passi pi\`{u} piccoli dove il moto evolve pi\`{u} rapidamente, es. vicino a un ostacolo). Questo \`{e} esattamente il tipo di scelta che un motore di simulazione fisica (es. per VR) deve fare a ogni frame.
\end{nota}

\begin{definizione}[Accelerazione Scalare]
\[
a_m = \bar a = \frac{\Delta v}{\Delta t} \quad \text{(media)} \qquad\qquad a = \frac{dv}{dt} = \frac{d^2x}{dt^2} \quad \text{(istantanea)}
\]
Relazioni inverse (utili da ricordare): \textbf{velocit\`{a} = integrale dell'accelerazione}; \textbf{posizione = integrale della velocit\`{a}}.
\end{definizione}

\begin{teorema}[Moto Rettilineo Uniforme]
Velocit\`{a} costante ($a=0$): $x = x_0 + vt$ (\textbf{legge oraria}).
\end{teorema}

\begin{teorema}[Moto Rettilineo Uniformemente Accelerato]
Accelerazione costante:
\[
v = v_0 + at \qquad x = x_0 + v_0t + \frac12 at^2 \qquad v^2 = v_0^2 + 2a\Delta x
\]
La terza equazione si ottiene eliminando $t$ dalle prime due. Esempio classico: caduta libera sulla Terra con $a=g=9.81\ \text{m s}^{-2}$.
\end{teorema}

\begin{nota}[Moto vario e approssimazione a tratti]
Anche con accelerazione \textbf{variabile}, il moto pu\`{o} essere diviso in intervalli $\Delta t$ abbastanza piccoli da essere ben approssimati da $a=\text{costante}$ in ciascuno -- il principio alla base dell'\textbf{integrazione numerica} usata per programmare una simulazione fisica.
\end{nota}

\subsection{Cinematica in 2D/3D: Traiettoria, Versori, Proiettili}

In $\mathbb{R}^2$ (o $\mathbb{R}^3$) un punto \`{e} definito da 2 (o 3) coordinate; lo spostamento $\Delta\vec r = \vec r(t+\Delta t)-\vec r(t)$ si trova per differenza vettoriale.

\begin{definizione}[Versori Tangente e Normale]
Lungo una traiettoria curva, in ogni punto si definiscono due versori: $\hat u_T$, sempre \textbf{tangente} alla traiettoria (punta nel verso del moto), e $\hat u_N$, sempre \textbf{perpendicolare} alla traiettoria, puntando verso il centro di curvatura locale.
\[
\vec v = v\,\hat u_T \qquad\qquad \vec a = a_T\hat u_T + a_N \hat u_N
\]
La componente \textbf{tangenziale} $a_T$ cambia il \emph{modulo} di $\vec v$; la componente \textbf{normale} $a_N$ ne cambia la \emph{direzione}.
\end{definizione}

\begin{teorema}[Moto del Proiettile]
Composizione di un moto \textbf{rettilineo uniforme} lungo $x$ e \textbf{uniformemente accelerato} (da $g$) lungo $y$:
\[
x = x_0 + v_{0x}t \qquad\qquad y = y_0 + v_{0y}t + \frac12 gt^2
\]
Eliminando $t$, la traiettoria \`{e} una \textbf{parabola}. Per un lancio da terra con angolo $\theta$ e velocit\`{a} $v_0$, la \textbf{gittata} (range) \`{e}:
\[
R = -\frac{v_0^2\sin(2\theta)}{g}
\]
(con la convenzione $g=-9.81\ \text{m s}^{-2}$, asse $y$ verso l'alto, cosicch\'{e} $R$ risulti positivo).
\end{teorema}

\begin{center}
\begin{tikzpicture}
\begin{axis}[width=9cm, height=5.5cm, xlabel={$x$ (m)}, ylabel={$y$ (m)}, axis lines=left, font=\small, ymin=0]
\addplot[domain=0:17.9, samples=100, thick, blue] {x*tan(60) - 9.81*x^2/(2*20^2*cos(60)^2)};
\node[align=center] at (axis cs:9,4) {\footnotesize $\varphi_0=60^{\circ}$, senza attrito\\ gittata $\approx 17.9$ m};
\end{axis}
\end{tikzpicture}
\captionof{figure}{Traiettoria parabolica ideale di un proiettile lanciato a $60^{\circ}$. In presenza di attrito viscoso dell'aria la traiettoria diventa asimmetrica (discesa pi\`{u} ripida della salita) e la gittata si riduce sensibilmente.}
\end{center}

\begin{nota}[Attrito dell'aria: attrito viscoso]
Nel caso reale, l'attrito dell'aria rende anche il moto lungo $x$ decelerato: la decelerazione non \`{e} costante ma dipende dalla velocit\`{a} del corpo (\textbf{attrito viscoso}), rendendo la traiettoria asimmetrica (discesa pi\`{u} ripida della salita).
\end{nota}

\begin{esempio}[Moto Relativo -- Composizione di Velocit\`{a}]
Se un lancio verticale avviene su uno skateboard in moto: nel riferimento dello skateboard il moto \`{e} solo lungo $y$ (uniformemente accelerato); nel riferimento dell'osservatore fermo, lo skateboard aggiunge una velocit\`{a} costante lungo $x$, dando luogo a un moto parabolico.
\[
\vec v_A = \vec v_R + \vec v_T
\]
dove $\vec v_A$ = velocit\`{a} \textbf{assoluta} (rispetto all'osservatore fermo), $\vec v_R$ = velocit\`{a} \textbf{relativa} (rispetto al sistema in moto), $\vec v_T$ = velocit\`{a} di \textbf{trascinamento} (del sistema in moto rispetto a quello fermo). Applicazione classica: una barca che attraversa un fiume con corrente, dove $\vec v_{BS} = \vec v_{BW} + \vec v_{WS}$ (barca-sponda = barca-acqua + acqua-sponda).
\end{esempio}

\subsection{Moto Circolare Uniforme}

\begin{definizione}[Moto Circolare Uniforme]
Moto lungo un cerchio con velocit\`{a} \textbf{costante in modulo}. \textbf{Attenzione}: essendo $\vec v$ un vettore, ``uniforme'' si riferisce solo al modulo -- la direzione cambia continuamente, quindi esiste comunque un'accelerazione. Il moto circolare uniforme \`{e} in realt\`{a} un moto bidimensionale uniformemente accelerato.
\end{definizione}

\begin{teorema}[Accelerazione Centripeta]
Per geometria (triangoli simili tra il triangolo delle velocit\`{a} e quello degli spostamenti):
\[
a = \frac{dv}{dt} = \frac{v^2}{r}
\]
detta \textbf{accelerazione centripeta} perch\'{e} punta sempre verso il centro.

\textbf{Grandezze fondamentali}: periodo $T$ [s]; frequenza $f=1/T$ [Hz]; frequenza angolare (pulsazione) $\omega = 2\pi/T = 2\pi f$ [rad/s]; velocit\`{a} $v = \omega r$.
\end{teorema}

%=============================================================================
\section{Dinamica}
%=============================================================================

\subsection{Forza e Principi della Dinamica}

\begin{definizione}[Dinamica e Forza]
La \textbf{dinamica} studia le \emph{cause} del moto: alla base c'\`{e} il concetto di \textbf{forza} come causa di una variazione di moto (o di deformazione). Le forze sono quantit\`{a} \textbf{vettoriali}, e si dividono in:
\begin{itemize}
  \item \textbf{Forze di contatto}: richiedono un punto di applicazione (es. una mano che tira una molla).
  \item \textbf{Forze di campo}: agiscono a distanza, permeando lo spazio (es. gravit\`{a}, elettrostatica, magnetismo).
\end{itemize}
In realt\`{a} le forze di contatto sono forze di campo (elettromagnetiche) a livello microscopico.
\end{definizione}

\begin{teorema}[Le Tre Leggi di Newton]
\begin{enumerate}
  \item \textbf{Principio d'inerzia}: in assenza di forze esterne, un corpo mantiene il suo stato di quiete o di moto rettilineo uniforme. Nei \textbf{sistemi inerziali} vale il principio d'inerzia; nei \textbf{sistemi non inerziali} (in moto accelerato) servono forze apparenti per spiegare variazioni di moto senza causa evidente.
  \item \textbf{Legge di Newton}: $\sum_i \vec F_i = m\vec a$, dove $m$ (massa inerziale, scalare) \`{e} la costante di proporzionalit\`{a} tra forza risultante e accelerazione. Unit\`{a} di forza: N $=$ kg m s$^{-2}$.
  \item \textbf{Azione e reazione}: $\vec F_{AB} = -\vec F_{BA}$ -- stesso modulo e direzione, verso opposto. \textbf{Le due forze agiscono su corpi diversi} (es. martello-chiodo: la forza del martello sul chiodo e quella del chiodo sul martello).
\end{enumerate}
\end{teorema}

\begin{teorema}[Dinamica del Moto Circolare Uniforme]
Riprendendo l'accelerazione centripeta (cinematica), la \textbf{forza centripeta} responsabile \`{e}:
\[
F_c = ma_c = m\frac{v^2}{r}
\]
Se il vincolo che fornisce questa forza viene meno (es. una corda si spezza), il corpo prosegue di moto rettilineo uniforme \textbf{tangente} al cerchio nel punto di rilascio (I legge) -- \textbf{non} continua a curvare.
\end{teorema}

\subsection{Forza Gravitazionale ed Elastica}

\begin{teorema}[Legge di Gravitazione Universale -- Newton, 1687]
\[
\vec F_{12} = -\frac{Gm_1m_2}{|\vec r_{12}|^2}\hat u_r \qquad G = 6.67\times10^{-11}\ \text{N m}^2\text{kg}^{-2}
\]
Forza attrattiva (segno negativo), a range infinito, proporzionale al prodotto delle masse. Accoppiando con la II legge di Newton per un corpo di massa $m$ sulla superficie terrestre:
\[
mg = \frac{GM_\oplus m}{R_\oplus^2} \quad\Rightarrow\quad g = \frac{GM_\oplus}{R_\oplus^2} \approx 9.81\ \text{m s}^{-2}
\]
con $M_\oplus = 5.972\times10^{24}$ kg, $R_\oplus = 6.371\times10^6$ m. \textbf{Importante}: $g$ non dipende dalla massa del corpo (massa inerziale = massa gravitazionale), ma \textbf{non} \`{e} costante: varia con l'altitudine (es. sull'Everest $g\approx 9.78\ \text{m s}^{-2}$, circa 0.3\% in meno). La \textbf{forza peso} $\vec P = m\vec g$ \`{e} proporzionale alla massa -- attenzione: \textbf{peso $\neq$ massa}.
\end{teorema}

\begin{teorema}[Legge di Hooke -- Forza Elastica]
\[
\vec F = -k(\vec x - \vec x_0)
\]
Forza di \textbf{richiamo}: sempre opposta allo spostamento, punta sempre verso la posizione di equilibrio (sia in compressione che in allungamento). Un disegno che indica tutte le forze agenti su un corpo si chiama \textbf{diagramma di corpo libero} (free-body diagram).
\end{teorema}

\begin{esempio}[La forza elastica nelle simulazioni 3D: spring-mass systems]
Nelle simulazioni fisiche per computer grafica/VR, i vertici di una mesh poligonale sono trattati come un sistema di masse connesse da molle (ciascun vertice connesso tipicamente a 1-4 vicini), di tre tipi: \textbf{structural spring} (struttura portante), \textbf{shear spring} (resistenza al taglio), \textbf{bend spring} (resistenza alla flessione). Questa combinazione \`{e} alla base della simulazione di tessuti (\emph{cloth simulation}, su una griglia regolare di vertici) e di corpi deformabili 3D (\emph{soft-body simulation}).
\end{esempio}

\subsection{Forza di Attrito e Vincoli}

\begin{definizione}[Forze Vincolari e Attrito]
I \textbf{vincoli} costringono il moto (es. in una direzione); generano una \textbf{reazione vincolare} $\vec F_N$ che bilancia la componente della forza risultante normale al vincolo (non c'\`{e} moto in quella direzione).

La \textbf{forza di attrito} (frenante, $F_{fr}$) si oppone alla forza applicata, con modulo proporzionale alla reazione normale:
\begin{itemize}
  \item \textbf{Statico} (nessun moto, $\sum F=0$): $F_{fr} \leq \mu_s F_N$ ($\mu_s$ = coefficiente di attrito statico).
  \item \textbf{Dinamico/cinetico} (c'\`{e} moto): $F_{fr} = \mu_d F_N$, con $\mu_d < \mu_s$.
\end{itemize}
\end{definizione}

\begin{center}
\begin{tikzpicture}
\begin{axis}[width=8cm, height=5.5cm, xlabel={Forza applicata $F_A$}, ylabel={Forza d'attrito $F_{fr}$}, xtick=\empty, ytick=\empty, axis lines=left, font=\small, ymax=45]
\addplot[domain=0:38, samples=50, thick, blue] {x};
\addplot[domain=38:70, samples=50, thick, blue] {30 + 1.2*sin(15*x)};
\draw[dashed] (axis cs:38,0) -- (axis cs:38,38);
\node[below] at (axis cs:38,0) {\footnotesize $\mu_sF_N$};
\node[align=center] at (axis cs:18,30) {\footnotesize nessun moto\\(attrito statico)};
\node[align=center] at (axis cs:55,20) {\footnotesize scivolamento\\(attrito dinamico)};
\end{axis}
\end{tikzpicture}
\captionof{figure}{L'attrito statico cresce linearmente con la forza applicata fino al valore massimo $\mu_sF_N$; superata questa soglia, il corpo scivola e l'attrito scende al valore dinamico $\mu_dF_N < \mu_sF_N$, oscillando durante lo scivolamento.}
\end{center}

\begin{esempio}[Piano Inclinato]
Su un corpo su un piano inclinato di angolo $\theta$ agiscono: la forza peso $\vec F_G=m\vec g$ (scomposta in $F_{Gx}=mg\sin\theta$ lungo il piano e $F_{Gy}=mg\cos\theta$ normale al piano), la reazione vincolare $F_N$ (che bilancia $F_{Gy}$), e l'attrito $F_{fr}=\mu F_N$ (lungo il piano, opposto al moto). La componente $F_{Gy}$, pur bilanciata dal vincolo, resta comunque necessaria per calcolare $F_{fr}$.
\end{esempio}

\begin{center}
\begin{tikzpicture}[font=\small]
\draw[thick] (0,0) -- (5,0) -- (5,2.2) -- cycle;
\draw (0.9,0) arc (0:23.7:0.9);
\node at (1.35,0.22) {\footnotesize $\theta$};
\coordinate (P) at (3,1.32);
\draw[fill=blue!15, thick] (P) ++(-15:-0.35) rectangle ++(0.5,0.5);
\draw[-Stealth, thick, red] (P) -- ++(90-24.6:1.1) node[above]{$F_N$};
\draw[-Stealth, thick, orange!80!black] (P) -- ++(180-24.6:0.9) node[left]{$F_{fr}$};
\draw[-Stealth, thick, blue] (P) -- ++(0,-1.3) node[below]{$F_G=mg$};
\draw[-Stealth, thick, blue, dashed] (P) -- ++(-24.6:0.9) node[below right]{$F_{Gx}$};
\draw[-Stealth, thick, blue, dashed] (P) -- ++(-114.6:0.7) node[right]{$F_{Gy}$};
\end{tikzpicture}
\captionof{figure}{Analisi delle forze su un piano inclinato: la forza peso si scompone in una componente lungo il piano ($F_{Gx}=mg\sin\theta$, che tende a far scivolare il corpo) e una normale al piano ($F_{Gy}=mg\cos\theta$, bilanciata da $F_N$).}
\end{center}

\begin{esempio}[Moto Circolare con Attrito -- Auto in Curva]
In un sistema non inerziale (l'auto in curva), la forza di attrito tra pneumatici e strada fornisce la forza centripeta necessaria. Bilancio verticale: $F_N = mg$. Bilancio radiale: $F_{fr} \leq \mu_s F_N = \mu_s mg = m v_R^2/r$, da cui la \textbf{velocit\`{a} massima in curva}:
\[
v_{R,\max} = \sqrt{\mu_s g r}
\]
\textbf{Domanda concettuale}: la forza percepita dai passeggeri \`{e} centripeta o centrifuga? Per la I legge, i passeggeri tendono a proseguire in linea retta (da cui la sensazione di forza ``centrifuga''); ma per la II legge, \`{e} in realt\`{a} una forza centripeta reale che li costringe sulla traiettoria curva -- la forza centrifuga \`{e} solo apparente.
\end{esempio}

\subsection{Le Tre Leggi di Keplero}

\begin{teorema}[Prima Legge di Keplero -- Orbite Ellittiche]
Tutti i pianeti si muovono su orbite \textbf{ellittiche} di cui il Sole occupa uno dei due fuochi. L'\textbf{eccentricit\`{a}} $e$ misura quanto l'ellisse \`{e} ``schiacciata'' ($e=0$: cerchio). Le orbite dei pianeti del sistema solare sono quasi circolari, eccetto Mercurio ($e=0.206$) e il pianeta nano Plutone ($e=0.249$).
\end{teorema}

\begin{teorema}[Seconda Legge di Keplero -- Legge delle Aree]
Il raggio vettore che congiunge un pianeta al Sole spazza aree uguali in tempi uguali (\textbf{velocit\`{a} aereolare costante}). Conseguenza (conservazione del momento angolare, cfr. rotazioni): un pianeta si muove \textbf{pi\`{u} velocemente al perielio} (punto pi\`{u} vicino al Sole) e \textbf{pi\`{u} lentamente all'afelio} (punto pi\`{u} lontano).
\end{teorema}

\begin{teorema}[Terza Legge di Keplero]
Il quadrato del periodo orbitale \`{e} proporzionale al cubo del semiasse maggiore dell'orbita. Per orbite quasi circolari ($R\sim a$, $M\gg$ massa del pianeta):
\[
T^2 = \frac{4\pi^2}{GM}R^3
\]
\end{teorema}

\begin{esempio}[Orbita Geostazionaria]
Un satellite geostazionario ha $T=86\,400$ s (un giorno). Risolvendo per $R$ con $M=M_\oplus=5.97\times10^{24}$ kg:
\[
R = \sqrt[3]{\frac{GMT^2}{4\pi^2}} = 4.223\times10^7\ \text{m} = 42\,230\ \text{km}
\]
Altitudine sul livello del mare: $h = R - R_\oplus = 42\,230-6\,371 = \mathbf{35\,900\ \text{km}}$.
\end{esempio}

%=============================================================================
\section{Lavoro, Energia e Leggi di Conservazione}
%=============================================================================

\subsection{Lavoro}

\begin{definizione}[Lavoro]
Prodotto \textbf{scalare} tra forza e spostamento:
\[
W = \vec F \cdot \vec d = Fd\cos\theta \qquad \text{Unit\`{a}: N m} = \text{J (Joule)}
\]
Se $\vec F \perp \vec d$ allora $W=0$ (es. una persona che trasporta una borsa in orizzontale: la forza verso l'alto non compie lavoro sullo spostamento orizzontale).
\end{definizione}

\begin{teorema}[Lavoro di una Forza Variabile]
Per una forza non costante, si suddivide il percorso in intervalli infinitesimi $dx$ su cui $F(x)$ \`{e} approssimativamente costante:
\[
W = \int_{d_A}^{d_B} \vec F(x)\cdot d\vec x
\]
Geometricamente, $W$ \`{e} l'\textbf{area sotto il grafico} $F$ vs. $x$.
\end{teorema}

\begin{center}
\begin{tikzpicture}
\begin{axis}[width=8.5cm, height=5.5cm, xlabel={$x$}, ylabel={$F_\parallel$}, xtick=\empty, ytick=\empty, axis lines=left, font=\small, ymin=0]
\addplot[domain=0.5:4, samples=100, thick, blue, name path=f] {30+15*x};
\path[name path=xaxis] (axis cs:0.5,0) -- (axis cs:4,0);
\addplot[domain=0.5:4, thick, black] coordinates {(0.5,0) (0.5,37.5)};
\addplot[domain=0.5:4, thick, black] coordinates {(4,0) (4,90)};
\addplot[blue!15] fill between[of=f and xaxis];
\node[below] at (axis cs:0.5,0) {$d_A$};
\node[below] at (axis cs:4,0) {$d_B$};
\node[align=center] at (axis cs:2.2,20) {\footnotesize $W=\displaystyle\int_{d_A}^{d_B}\! F(x)\,dx$};
\end{axis}
\end{tikzpicture}
\captionof{figure}{Il lavoro di una forza variabile \`{e} l'area sotto la curva $F(x)$ tra le posizioni iniziale e finale.}
\end{center}

\subsection{Energia Cinetica e Teorema delle Forze Vive}

\begin{definizione}[Energia ed Energia Cinetica]
L'\textbf{energia} \`{e} la capacit\`{a} di compiere lavoro. L'\textbf{energia cinetica} \`{e} la capacit\`{a} di compiere lavoro dovuta allo stato di moto:
\[
K = \frac12 mv^2 \qquad \text{(sempre} \geq 0\text{)}
\]
\end{definizione}

\begin{teorema}[Teorema delle Forze Vive (teorema dell'energia cinetica)]
\[
W = K_2 - K_1 = \Delta K
\]
Il lavoro compiuto su un corpo equivale alla variazione della sua energia cinetica. \textbf{Attenzione al segno}: $W>0$ se $K$ aumenta.
\end{teorema}

\begin{nota}[Da cosa, su cosa: il segno del lavoro]
Per la III legge, se un martello esercita forza $\vec F$ su un chiodo, il chiodo esercita $-\vec F$ sul martello: $W_{\text{sul chiodo}} = -W_{\text{sul martello}}$. Il martello \emph{perde} energia cinetica, il chiodo la \emph{guadagna}. Il ragionamento \`{e} intuitivo per coppie di oggetti, ma diventa pi\`{u} sottile quando si parla di lavoro compiuto da una singola \emph{forza} (es. gravit\`{a}) senza un oggetto reciproco esplicito.
\end{nota}

\subsection{Forze Conservative ed Energia Potenziale}

\begin{definizione}[Forze Conservative]
Una forza \`{e} \textbf{conservativa} se il lavoro da $A$ a $B$ non dipende dal percorso ($W_{AB}=W_{AC}+W_{CB}$), equivalentemente se il lavoro lungo un percorso \textbf{chiuso} \`{e} sempre nullo. Esempi: gravitazionale, elastica. Le forze \textbf{non conservative} (es. attrito) dipendono dal percorso e dissipano parte dell'energia in calore.
\end{definizione}

\begin{definizione}[Energia Potenziale]
L'energia potenziale \`{e} la capacit\`{a} di compiere lavoro dovuta alla \textbf{posizione} in un campo di forze conservative, definita \textbf{a meno di una costante arbitraria} $U_0$ (es. una quota di riferimento):
\[
U_{grav}(h) = mgh + U_0 \qquad\qquad U_{el}(x) = \frac12 kx^2 + U_0
\]
Relazione generale con il lavoro: $W = -\Delta U = U_i - U_f$ (il lavoro della forza conservativa \`{e} l'opposto della variazione di potenziale).
\end{definizione}

\begin{nota}[Costruzione intuitiva del potenziale gravitazionale]
Sollevando lentamente ($a=0$) un mattone di un'altezza $h$: la forza applicata bilancia il peso ($F_A=mg$), compiendo lavoro $W_A=mgh$ \emph{contro} la gravit\`{a} ($W_G=-mgh$). Lasciandolo cadere, la gravit\`{a} compie lavoro positivo $W_G=mgh$, e per il teorema delle forze vive il mattone acquista $K=mgh>0$: ha quindi acquisito la \emph{capacit\`{a}} di compiere altrettanto lavoro (es. conficcando un cuneo nel terreno all'impatto) -- da cui la definizione $U=mgh+U_0$.
\end{nota}

\begin{teorema}[Conservazione dell'Energia Meccanica]
\[
E_{tot} = K + U = \text{costante}
\]
vale per un \textbf{sistema isolato} (nessuna interazione esterna) in cui ogni corpo interagisce solo tramite forze \textbf{conservative} ($U=\sum U$, somma di tutti i potenziali presenti). L'energia meccanica totale \`{e} un ``integrale primo'' del moto: non dipende dal tempo n\'{e} dagli stati intermedi, solo dalle posizioni iniziale e finale.
\end{teorema}

\begin{center}
\begin{tikzpicture}[scale=0.9]
\draw[thick, blue] plot[smooth] coordinates {(0,2.6) (1.2,0.3) (2.4,2.6)};
\foreach \x/\lbl in {0/1, 1.2/2, 2.4/3} {
  \fill[red] (\x,{\x==1.2?0.3:2.6}) circle (2pt);
}
\node[above] at (0,2.6) {1};
\node[below] at (1.2,0.3) {2};
\node[above] at (2.4,2.6) {3};
\end{tikzpicture}
\captionof{figure}{Su una traiettoria a U (senza attrito), l'energia cinetica guadagnata cadendo da 1 a 2 viene interamente riconvertita in energia potenziale risalendo da 2 a 3 (stessa quota di 1): $\Delta U = -\Delta K$ vale per ogni coppia di punti, incluso 1$\to$3 direttamente.}
\end{center}

\begin{teorema}[Relazione tra Forza e Potenziale]
In una dimensione: $F_x = -\dfrac{dU}{dx}$ ($U$ \`{e} la primitiva di $-F$). In 3D, tramite il \textbf{gradiente}:
\[
\vec F = -\left(\frac{\partial U}{\partial x}\hat u_x+\frac{\partial U}{\partial y}\hat u_y+\frac{\partial U}{\partial z}\hat u_z\right) = -\vec\nabla U
\]
Le \textbf{superfici equipotenziali} (punti a $U$ costante) sono sempre \textbf{perpendicolari} alla forza in ogni loro punto (es. le curve di livello di una montagna, o le superfici sferiche concentriche attorno alla Terra su scala planetaria). Un'orbita \textbf{circolare} avviene su un'unica superficie equipotenziale ($U$ e $K$ separatamente costanti); un'orbita \textbf{ellittica} attraversa superfici equipotenziali diverse ($U+K$ costante, ma $U,K$ variano singolarmente -- pi\`{u} veloce vicino alla stella, per la II legge di Keplero).
\end{teorema}

\begin{nota}[Lavoro con forze non conservative]
Il teorema delle forze vive $W=\Delta K$ vale sempre. Il lavoro totale include sia il contributo conservativo che quello dissipativo: $W = W_{cons}+W_{diss}$, e l'energia totale cambia in modo calcolabile: $E_{tot,nuova} = E_{tot,vecchia} + W_{diss}$.
\end{nota}

\subsection{Potenza}

\begin{definizione}[Potenza]
Descrive quanto rapidamente si compie un lavoro:
\[
P = \frac{W}{\Delta t}\ \text{(media)} \qquad\qquad P = \frac{dW}{dt} = \vec F\cdot\vec v\ \text{(istantanea, per } F,v \text{ costanti)}
\]
Unit\`{a}: W (Watt) $=$ J/s. A differenza dell'energia potenziale, il lavoro $W$ non ha un valore di riferimento arbitrario da sottrarre: \`{e} una quantit\`{a} assoluta per un dato processo.
\end{definizione}

\subsection{Fisica ed Energia in un Motore Fisico}

\begin{nota}[Schema di un motore fisico (physics engine)]
Un ciclo di simulazione tipico, a intervalli temporali $dt$ fissati:
\begin{enumerate}
  \item Inizializzare $x,v,a$ per ogni punto materiale di massa $m$.
  \item Per ogni ciclo di durata $dt$: risolvere le forze ($\vec a = \vec F/m$); aggiornare $\vec v_{new}=\vec v_{old}+\vec a\,dt$; aggiornare $\vec r_{new}=\vec r_{old}+\vec v\,dt$ (integrazione di Eulero).
  \item Identificare e risolvere eventuali collisioni.
\end{enumerate}
Questo \`{e} lo scheletro base di un motore fisico in tempo reale (es. per VR/videogiochi).
\end{nota}

\begin{nota}[La conservazione dell'energia come controllo di correttezza]
In un motore fisico, la conservazione dell'energia meccanica pu\`{o} essere usata come \textbf{verifica} che le approssimazioni numeriche siano accettabili:
\begin{itemize}
  \item Tutte le quantit\`{a} sono rappresentate con numeri in virgola mobile (\emph{float}), con errore numerico finito.
  \item Intervalli $dt$ troppo grandi rendono l'approssimazione $a=\text{costante}$ meno accurata.
  \item I metodi di integrazione numerica tendono a far \textbf{aumentare o diminuire} artificialmente l'energia nel tempo (\emph{energy drift}).
\end{itemize}
Ad ogni ciclo si pu\`{o} verificare (entro una soglia d'errore accettabile) che $E = K+U_{grav}+U_{el}-E_{dissipata}$ resti costante e che $W=\Delta K$ (teorema delle forze vive); se il controllo fallisce, si applicano correzioni.
\end{nota}

%=============================================================================
\section{Oscillatore Armonico e Pendolo}
%=============================================================================

\subsection{Equilibrio ed Energia Potenziale}

\begin{definizione}[Tipi di Equilibrio]
Data una curva di energia potenziale $U(x)$, con $\vec F = -dU/dx$:
\begin{itemize}
  \item \textbf{Equilibrio stabile}: minimo di $U$; una perturbazione genera una \textbf{forza di richiamo}.
  \item \textbf{Equilibrio instabile}: massimo di $U$; una perturbazione genera una \textbf{forza repulsiva}.
  \item \textbf{Equilibrio indifferente}: $U$ costante; una perturbazione non genera forze.
\end{itemize}
\end{definizione}

\subsection{L'Oscillatore Armonico}

\begin{definizione}[Oscillatore Armonico]
Sistema in cui la legge di Hooke ($F=-kx$) causa un moto oscillatorio periodico attorno a un equilibrio stabile, con potenziale parabolico $U=\frac12 kx^2$. L'\textbf{ampiezza} $A$ \`{e} l'elongazione massima ($-A\leq x\leq A$).
\end{definizione}

\begin{teorema}[Energia dell'Oscillatore Armonico]
La forza elastica \`{e} conservativa:
\[
E_{tot} = K+U_{el} = \frac12 mv^2+\frac12 kx^2 = \text{costante} = \frac12 kA^2
\]
(valutando a $x=A$, dove $v=0$). \textbf{L'energia meccanica \`{e} proporzionale al quadrato dell'ampiezza.} Da qui si ricava la velocit\`{a} a elongazione generica:
\[
v = \pm v_0\sqrt{1-\frac{x^2}{A^2}} \qquad v_0 = \sqrt{\frac{k}{m}}A \ \text{(velocit\`{a} massima, a } x=0\text{)}
\]
\end{teorema}

\begin{nota}[Moto armonico come proiezione del moto circolare]
Un punto che si muove di moto circolare uniforme con raggio $A$ e velocit\`{a} orbitale $v_0$, proiettato su un asse (es. $x$), ha $x=A\cos\theta$ e velocit\`{a} proiettata $v=v_0\sin\theta = v_0\sqrt{1-x^2/A^2}$ -- esattamente la stessa relazione ricavata dalla conservazione dell'energia. Questa analogia \`{e} il modo pi\`{u} intuitivo per derivare periodo, legge oraria, velocit\`{a} e accelerazione del moto armonico.
\end{nota}

\begin{teorema}[Periodo, Frequenza e Legge Oraria]
Dall'analogia con il moto circolare ($T=2\pi A/v_0$) e $v_0=\sqrt{k/m}\,A$:
\[
T = 2\pi\sqrt{\frac{m}{k}} \qquad f=\frac1T=\frac{1}{2\pi}\sqrt{\frac{k}{m}} \qquad \omega = 2\pi f = \sqrt{\frac{k}{m}} \ \text{(pulsazione)}
\]
\textbf{Il periodo non dipende dall'ampiezza.} La legge oraria (posizione, velocit\`{a}, accelerazione):
\[
x(t)=A\cos(\omega t+\varphi) \qquad v(t)=-\omega A\sin(\omega t+\varphi) \qquad a(t)=-\omega^2A\cos(\omega t+\varphi)=-\omega^2x(t)
\]
La velocit\`{a} \`{e} sfasata di $90^{\circ}$ rispetto a $x$ (massima all'equilibrio); l'accelerazione \`{e} in fase con $-x$ (sempre di richiamo, come la forza elastica).
\end{teorema}

\begin{center}
\begin{tikzpicture}
\begin{axis}[width=9.5cm, height=6cm, xlabel={$t$}, xtick=\empty, ytick=\empty, axis lines=middle, font=\small, ymin=-1.5, ymax=1.5, axis line style={-Stealth}]
\addplot[domain=0:4*pi, samples=150, thick, blue] {cos(deg(x))};
\addplot[domain=0:4*pi, samples=150, thick, red!70!black] {-sin(deg(x))};
\addplot[domain=0:4*pi, samples=150, thick, green!50!black] {-cos(deg(x))};
\legend{$x(t)/A$, $v(t)/v_0$, $a(t)\cdot m/(kA)$}
\end{axis}
\end{tikzpicture}
\captionof{figure}{Posizione, velocit\`{a} e accelerazione nel moto armonico: la velocit\`{a} \`{e} sfasata di $90^{\circ}$ rispetto alla posizione, l'accelerazione \`{e} in opposizione di fase (sempre diretta verso l'equilibrio).}
\end{center}

\begin{esempio}[Molla Verticale in Presenza di Gravit\`{a}]
Appendendo la molla in verticale, il nuovo equilibrio si sposta a $x_0=mg/k$ (dove $mg=kx_0$), ma il moto attorno a questa nuova posizione \`{e} identico al caso orizzontale: $T,f,\omega,k$ restano invariati -- solo il punto di equilibrio cambia.
\end{esempio}

\begin{nota}[Cenni al moto armonico smorzato]
Aggiungendo una forza di attrito viscoso proporzionale alla velocit\`{a} ($\vec F_{fr}=-b\vec v$), l'equazione del moto diventa un'\textbf{equazione differenziale}:
\[
m\frac{d^2x}{dt^2}+b\frac{dx}{dt}+kx=0
\]
con soluzione $x(t)=A\cos(\omega t+\varphi)\,e^{-t(b/2m)}$: un'oscillazione la cui ampiezza decresce esponenzialmente nel tempo.
\end{nota}

\begin{nota}[Perch\'{e} l'oscillatore armonico \`{e} cos\`{i} importante]
Qualunque potenziale, anche complesso, pu\`{o} essere approssimato da una parabola nei pressi di un suo minimo (sviluppo in serie di Taylor troncato al secondo ordine). Per questo l'oscillatore armonico descrive bene il comportamento di sistemi anche complessi vicino a un equilibrio stabile -- dagli atomi di idrogeno alle perturbazioni di un oggetto in un videogioco o in una simulazione VR.
\end{nota}

\subsection{Il Pendolo Semplice}

\begin{definizione}[Pendolo Semplice]
Filo inestensibile e senza massa, con un punto materiale $m$ appeso all'estremit\`{a}: un moto armonico sotto l'effetto della gravit\`{a}.
\end{definizione}

\begin{teorema}[Equazione del Pendolo e Approssimazione per Piccole Oscillazioni]
Scomponendo la forza peso in componente radiale ($mg\cos\theta$, bilanciata dalla tensione $T$) e tangenziale ($-mg\sin\theta$, di richiamo), e usando $s=L\theta$ (arco):
\[
\frac{d^2\theta}{dt^2} = -\frac{g}{L}\sin\theta \ \xrightarrow{\ \theta\lesssim 10^{\circ},\ \sin\theta\approx\theta\ }\ \frac{d^2\theta}{dt^2} = -\frac{g}{L}\theta
\]
Struttura identica all'oscillatore armonico ($a(t)=-\omega^2x(t)$, con $\theta$ al posto di $x$):
\[
\theta(t) = \theta_0\cos(\omega t) \qquad \omega=\sqrt{\frac{g}{L}} \qquad T = \frac{2\pi}{\omega}=2\pi\sqrt{\frac{L}{g}}
\]
Il periodo dipende solo da $g$ ed $L$ (\textbf{non dalla massa}), e non dipende dall'ampiezza finch\'{e} $\theta\lesssim10^{\circ}$ (\textbf{approssimazione delle piccole oscillazioni}). Per piccole oscillazioni, il pendolo si comporta come una molla dove $mg/L$ sostituisce $k$ (stesse unit\`{a}, N/m).
\end{teorema}

\begin{nota}[Attenzione alla notazione: $T$]
In questa sezione $T$ \`{e} usato sia per la \textbf{tensione} del filo (forza, in N) sia per il \textbf{periodo} di oscillazione (tempo, in s) -- due grandezze completamente diverse che condividono per convenzione lo stesso simbolo. Il contesto (un'equazione di forze vs. una legge oraria) chiarisce sempre quale delle due si intende.
\end{nota}

%=============================================================================
\section{Meccanica dei Sistemi: Collisioni e Traslazioni}
%=============================================================================

\subsection{Quantit\`{a} di Moto e Impulso}

\begin{nota}[Perch\'{e} servono nuovi strumenti]
Alcuni problemi (in particolare le \textbf{collisioni}) sono difficili da risolvere direttamente con $\sum\vec F=m\vec a$, perch\'{e} le forze di interazione interne durante l'urto sono complesse e spesso sconosciute. La \textbf{quantit\`{a} di moto} e la sua conservazione permettono di aggirare il problema.
\end{nota}

\begin{definizione}[Quantit\`{a} di Moto e II Legge Generalizzata]
\[
\vec p = m\vec v \qquad \text{Unit\`{a}: kg m s}^{-1}
\]
\[
\sum\vec F = \frac{d\vec p}{dt} = \frac{dm}{dt}\vec v + m\vec a
\]
Per massa costante si recupera $\sum\vec F=m\vec a$; la forma generalizzata copre anche sistemi a massa variabile.
\end{definizione}

\begin{teorema}[Conservazione della Quantit\`{a} di Moto]
Se $\sum\vec F=0$, allora $\vec p=\text{costante}$ (legge d'inerzia generalizzata). \textbf{Fondamentale per le collisioni}: se le forze esterne sono trascurabili rispetto a quelle interne durante l'urto, la quantit\`{a} di moto totale del sistema si conserva \emph{sempre}, indipendentemente dalla natura (nota o ignota) delle forze interne -- per la III legge, la somma delle forze interne \`{e} comunque nulla. \textbf{Attenzione}: essendo $\vec p$ un vettore, la conservazione vale componente per componente.
\end{teorema}

\begin{definizione}[Impulso]
\[
\vec I = \sum\vec F\,\Delta t\ \text{(forza costante)} \qquad\qquad \vec I = \int_{t_0}^{t_1}\vec F\,dt\ \text{(forza variabile)}
\]
\end{definizione}

\begin{teorema}[Teorema dell'Impulso]
\[
\vec I = \Delta\vec p
\]
L'impulso \`{e} pari alla variazione della quantit\`{a} di moto -- l'analogo, per la quantit\`{a} di moto, di quello che il teorema delle forze vive \`{e} per l'energia cinetica:
\begin{center}
\small
\begin{tabular}{ll}
\toprule
\textbf{Energia (scalare)} & \textbf{Quantit\`{a} di moto (vettore)} \\
\midrule
$K=\frac12mv^2$ & $\vec p = m\vec v$ \\
$W=\Delta K$ (lavoro $\leftrightarrow$ \emph{distanza}) & $\vec I=\Delta\vec p$ (impulso $\leftrightarrow$ \emph{tempo}) \\
\bottomrule
\end{tabular}
\end{center}
\end{teorema}

\begin{center}
\begin{tikzpicture}
\begin{axis}[width=8cm, height=5cm, xlabel={$t$}, ylabel={$F$}, xtick=\empty, ytick=\empty, axis lines=left, font=\small, ymin=0]
\addplot[domain=0:6, samples=100, thick, blue, name path=f] {40*exp(-((x-3)^2)/2)};
\path[name path=xaxis] (axis cs:0,0) -- (axis cs:6,0);
\addplot[thick, dashed, red] coordinates {(0,26) (6,26)};
\addplot[blue!12] fill between[of=f and xaxis];
\node[anchor=west, red] at (axis cs:6,26) {$\langle F\rangle$};
\node[below] at (axis cs:0,0) {$t_0$};
\node[below] at (axis cs:6,0) {$t_1$};
\end{axis}
\end{tikzpicture}
\captionof{figure}{La forza reale durante una collisione (curva blu) varia rapidamente nel tempo; l'area sotto la curva (impulso reale) \`{e} uguale all'area del rettangolo di altezza $\langle F\rangle$ (forza media) e base $\Delta t=t_1-t_0$ -- utile per stimare l'impulso senza conoscere il dettaglio di $F(t)$.}
\end{center}

\begin{nota}[L'impulso e la sua importanza pratica]
Quanto pi\`{u} a lungo si applica un'azione (una forza), tanto pi\`{u} cambia lo stato di moto dell'oggetto perturbato ($\vec I=\vec F\Delta t=\Delta\vec p$). Un oggetto di massa enorme richiede o una forza enorme per poco tempo, o una piccola forza applicata a lungo (es. deviare un asteroide con una piccola spinta sostenuta per anni, invece che con un'esplosione istantanea). Nel contesto di un'interfaccia VR, questo \`{e} lo stesso principio per cui l'effetto di premere il grilletto di un controller dipende da quanto a lungo lo si tiene premuto.
\end{nota}

\subsection{Collisioni}

\begin{definizione}[Collisione]
Un'interazione (anche complessa, con possibile perdita di energia) tra due o pi\`{u} corpi, in cui tipicamente le forze interne durante l'urto sono $\gg$ delle forze esterne, e quindi la quantit\`{a} di moto totale si conserva sempre:
\[
m_A\vec v_{A1}+m_B\vec v_{B1} = m_A\vec v_{A2}+m_B\vec v_{B2}
\]
\end{definizione}

\begin{teorema}[Collisione Completamente Anelastica]
I due corpi restano uniti dopo l'urto, con velocit\`{a} finale comune:
\[
\vec v_2 = \frac{m_A\vec v_{A1}+m_B\vec v_{B1}}{m_A+m_B}
\]
L'energia meccanica \textbf{non} si conserva ($K_{finale}<K_{iniziale}$): parte \`{e} dissipata nella deformazione/calore dell'urto. \textbf{Attenzione}: nel sistema di riferimento del centro di massa, \emph{tutta} l'energia cinetica relativa si perde; nel sistema del laboratorio, resta comunque l'energia cinetica del centro di massa (che si muove a $v_{CM}\neq0$).
\end{teorema}

\begin{teorema}[Collisione Elastica]
Sia la quantit\`{a} di moto sia l'energia cinetica si conservano (forze interne conservative). Risolvendo il sistema per un bersaglio inizialmente fermo ($v_{B1}=0$):
\[
v_{A2} = v_{A1}\,\frac{m_A-m_B}{m_A+m_B} \qquad\qquad v_{B2} = v_{A1}\,\frac{2m_A}{m_A+m_B}
\]
\textbf{Casi limite}:
\begin{itemize}
  \item $m_A\gg m_B$: $v_{A2}\approx v_{A1}$, $v_{B2}\approx2v_{A1}$ (una palla da bowling urta una da ping-pong: rallenta appena, la palla leggera parte al doppio della velocit\`{a}).
  \item $m_A\ll m_B$: $v_{A2}\approx-v_{A1}$, $v_{B2}\ll v_{A1}$ (ping-pong su bowling: la palla leggera rimbalza indietro quasi alla stessa velocit\`{a}, quella pesante si muove appena).
  \item $m_A=m_B$: $v_{A2}=0$, $v_{B2}=v_{A1}$ (due palle da biliardo: trasferimento completo di velocit\`{a}).
\end{itemize}
\end{teorema}

\begin{nota}[Collisioni nel sistema di riferimento del centro di massa]
Nel sistema del CM (dove $\vec v_{CM}=0$ per definizione), le quantit\`{a} di moto di due corpi che collidono sono sempre \textbf{uguali e opposte}, sia prima che dopo l'urto. Nell'urto \textbf{elastico}, ciascun momento si inverte semplicemente di segno mantenendo il modulo; nell'urto \textbf{completamente anelastico}, dopo l'urto entrambi i momenti si annullano (i corpi, uniti, coincidono con il CM, fermo in questo sistema di riferimento).
\end{nota}

\subsection{Il Centro di Massa}

\begin{nota}[Dal punto materiale al corpo esteso]
Abbandoniamo l'idealizzazione di ``punto materiale'' per descrivere sistemi di punti (discreti) o corpi estesi (continui). Vogliamo una legge generalizzata, tipo $\sum\vec F=m\vec a$, valida per le \textbf{traslazioni} di un intero sistema.
\end{nota}

\begin{definizione}[Centro di Massa]
La media ``pesata'' (con peso la massa) delle posizioni di tutti i punti del sistema:
\[
x_{CM} = \frac{1}{M}\sum_i m_i x_i \qquad M=\sum_i m_i \text{ (massa totale)}
\]
\end{definizione}

\begin{teorema}[Propriet\`{a} Fondamentale del Centro di Massa]
Concentrando tutta la massa del sistema nel CM e applicando $\sum\vec F_{est}=M\vec a_{CM}$, il \textbf{CM descrive il moto traslatorio dell'intero sistema} come se fosse un singolo punto materiale, indipendentemente dalle forze interne (che si annullano a coppie per la III legge di Newton). Questo vale \emph{solo} per il moto di traslazione: il CM non descrive le rotazioni del sistema attorno a s\'{e} stesso.

\emph{Dimostrazione (sistema discreto)}: sommando la II legge per ogni punto $i$, $\vec F_i+\sum_{j\neq i}\vec F_{ij}=m_i\vec a_i$, su tutti gli $i$; le forze interne $\vec F_{ij}$ si cancellano a coppie (III legge: $\vec F_{ij}=-\vec F_{ji}$), lasciando $\sum_i\vec F_i = \sum_i m_i\vec a_i = M\vec a_{CM}$.
\end{teorema}

\begin{esempio}[Tuffo con Rotazione]
Un tuffatore che ruota (capriole) in aria segue comunque, con il proprio centro di massa, l'identica traiettoria parabolica che seguirebbe un semplice punto materiale soggetto solo alla gravit\`{a} -- indipendentemente da come il corpo ruota o cambia posa (raccolto/disteso) durante il volo. Questo illustra visivamente che il CM descrive \emph{solo} la traslazione, non la rotazione.
\end{esempio}

\begin{nota}[Baricentro vs. Centro di Massa]
Il \textbf{baricentro} \`{e} il punto di applicazione della risultante della forza peso sui vari punti del sistema. Baricentro e centro di massa \textbf{coincidono} solo se il campo gravitazionale \`{e} uniforme (accelerazione di gravit\`{a} costante su tutto il corpo) -- condizione tipicamente ben verificata su scala umana/terrestre.
\end{nota}

\begin{esempio}[Il Centro di Massa nel Salto in Alto -- Tecnica Fosbury]
Durante il salto Fosbury, l'atleta inarca il corpo sopra l'asticella mentre il suo \textbf{centro di massa passa sotto} l'asticella. Poich\'{e} il lavoro necessario al salto dipende dall'innalzamento del CM ($W=Mgh_{CM}$, cfr. \S4), questa tecnica minimizza il lavoro richiesto per superare una data altezza dell'asticella, pur facendo passare l'intero corpo sopra di essa.
\end{esempio}

%=============================================================================
\section{Meccanica dei Sistemi: Moti di Rotazione}
%=============================================================================

\subsection{Cinematica Rotazionale}

\begin{definizione}[Corpo Rigido e Cinematica Angolare]
In un \textbf{corpo rigido}, ogni elemento ruota con la stessa velocit\`{a} angolare $\omega$ attorno a un asse di rotazione. Per un punto a distanza $r$ dall'asse: $v=\omega r$; se $\omega$ \`{e} costante il moto \`{e} circolare uniforme, con accelerazione puramente radiale (centripeta) $a_R=v^2/r=\omega^2r$.
\end{definizione}

\begin{teorema}[Accelerazione Totale nel Moto Circolare Vario]
Se $\omega$ \`{e} \textbf{variabile}, compare anche un'\textbf{accelerazione tangenziale}:
\[
a_T = \frac{dv}{dt} = r\ddot\theta \qquad\qquad \vec a = -\frac{v^2}{r}\hat u_r + r\ddot\theta\,\hat u_\theta
\]
(qui il punto sopra una variabile denota derivata temporale: $\dot\theta=\omega$, $\ddot\theta=\alpha$). Il versore radiale $\hat u_r$ punta verso l'esterno (da cui il segno meno davanti al termine centripeto).
\end{teorema}

\begin{teorema}[Analogia Traslazione-Rotazione]
Con accelerazione angolare $\alpha=\ddot\theta$ costante, le leggi rotazionali sono l'esatto analogo di quelle traslazionali ($\theta\leftrightarrow x$, $\omega\leftrightarrow v$, $\alpha\leftrightarrow a$):
\begin{center}
\small
\begin{tabular}{ll}
\toprule
\textbf{Rotazionale} & \textbf{Traslazionale} \\
\midrule
$\omega=\omega_0+\alpha t$ & $v=v_0+at$ \\
$\theta=\omega_0t+\frac12\alpha t^2$ & $s=v_0t+\frac12at^2$ \\
$\omega^2=\omega_0^2+2\alpha\theta$ & $v^2=v_0^2+2as$ \\
\bottomrule
\end{tabular}
\end{center}
\end{teorema}

\begin{definizione}[Puro Rotolamento]
Rotazione \textbf{senza} traslazione/slittamento relativo al suolo. Condizione: $v=\omega r$ (velocit\`{a} del centro rispetto al suolo = velocit\`{a} angolare per raggio). Visto dal sistema solidale con la ruota, il punto di contatto con il suolo si muove con velocit\`{a} $-v$.
\end{definizione}

\subsection{Momento Torcente e Momento d'Inerzia}

\begin{definizione}[Momento Torcente (Torque)]
Prodotto vettoriale tra il \textbf{braccio} (vettore dall'asse di rotazione al punto di applicazione) e la forza:
\[
\vec\tau = \vec r \wedge \vec F \qquad \tau = rF\sin\theta = r_\perp F = rF_\perp
\]
dove $r_\perp$ \`{e} la componente del braccio perpendicolare a $\vec F$, e $F_\perp$ la componente di $\vec F$ perpendicolare al braccio.
\end{definizione}

\begin{nota}[Forza vs. momento torcente]
Nel moto traslatorio, la forza \`{e} la causa dell'accelerazione lineare; bilanciare le forze evita un cambiamento nel moto di traslazione. Nel moto rotatorio, il \textbf{momento torcente} \`{e} la causa dell'accelerazione angolare; per evitare un cambio di rotazione occorre bilanciare i momenti torcenti, non le sole forze.
\end{nota}

\begin{esempio}[Vincoli e Momento Torcente -- il Volano]
Un volano su un asse verticale bloccato da vincoli agli estremi $A$ (basso) e $B$ (alto): l'asse non pu\`{o} n\'{e} traslare n\'{e} inclinarsi. Una forza $\vec F$ applicata al bordo si scompone in $\vec F_\perp$ (perpendicolare all'asse) e $\vec F_\parallel$ (parallela). Una reazione vincolare $R_\perp$ applicata \emph{sull'asse stesso} annulla la forza $F_\perp$ ma \textbf{non} il suo momento (braccio nullo per $R_\perp$, braccio pari al raggio per $F_\perp$) -- il volano pu\`{o} comunque ruotare. Analogamente, il momento di $F_\parallel$ non pu\`{o} essere bilanciato da un'unica reazione $R_\parallel$: serve una \textbf{coppia} di forze uguali e opposte $R',R''$ ai due estremi $A,B$ (risultante nulla, ma momento torcente $\neq 0$) per impedire l'inclinazione dell'asse.
\end{esempio}

\begin{teorema}[II Legge di Newton per le Rotazioni]
\[
\sum\vec\tau = I\vec\alpha
\]
L'accelerazione angolare \`{e} direttamente proporzionale al momento torcente risultante e inversamente proporzionale al \textbf{momento d'inerzia} $I$ (l'analogo rotazionale della massa inerziale). \textbf{Attenzione}: $\vec\omega$ e $\vec\alpha$ sono vettori, con verso dato dalla regola della mano destra (le dita seguono il verso di rotazione, il pollice indica $\vec\omega$).
\end{teorema}

\begin{definizione}[Momento d'Inerzia]
La resistenza di un corpo al cambiamento del proprio moto rotazionale:
\[
I = \sum_i m_id_i^2 \ \text{(discreto)} \qquad\qquad I = \int_V \rho\,d^2\,dV \ \text{(continuo)}
\]
dove $d_i$ (o $d$) \`{e} la distanza dall'asse di rotazione. \textbf{Il momento d'inerzia dipende dall'asse scelto}: a differenza della massa, non basta conoscere quanta materia c'\`{e}, ma anche \emph{come} \`{e} distribuita rispetto all'asse (un disco e un cilindro della stessa massa hanno momenti d'inerzia molto diversi a seconda dell'asse di rotazione).
\end{definizione}

\begin{center}
\small
\begin{tabular}{lll}
\toprule
\textbf{Oggetto} & \textbf{Asse} & \textbf{Momento d'inerzia} \\
\midrule
Anello sottile, raggio $R$ & centrale, perpendicolare al piano & $I=MR^2$ \\
Anello sottile, raggio $R$, spessore $W$ & diametro centrale & $I=\frac12MR^2+\frac{1}{12}MW^2$ \\
Cilindro (disco) pieno, raggio $R$ & centrale longitudinale & $I=\frac12MR^2$ \\
Cilindro cavo, raggi $R_1,R_2$ & centrale longitudinale & $I=\frac12M(R_1^2+R_2^2)$ \\
Sfera piena, raggio $R$ & centrale & $I=\frac25MR^2$ \\
Barra sottile, lunghezza $L$ & centrale, perpendicolare & $I=\frac{1}{12}ML^2$ \\
Barra sottile, lunghezza $L$ & un estremo, perpendicolare & $I=\frac13ML^2$ \\
Lamina rettangolare, $L\times W$ & centrale, perpendicolare & $I=\frac{1}{12}M(L^2+W^2)$ \\
\bottomrule
\end{tabular}
\captionof{table}{Momenti d'inerzia per solidi omogenei comuni (rispetto ad assi passanti per il centro geometrico, salvo dove indicato).}
\end{center}

\subsection{Energia, Rotolamento, Momento Angolare}

\begin{teorema}[Energia Cinetica Rotazionale e Totale]
Sommando l'energia cinetica di ciascun punto ($v_i=r_i\omega$) e raccogliendo $I=\sum m_ir_i^2$:
\[
K_{rot} = \frac12 I\omega^2
\]
Per un corpo che trasla e ruota insieme (es. rotolamento), rispetto al centro di massa (CM):
\[
K_{tot} = K_{tr}+K_{rot} = \frac12 Mv_{CM}^2 + \frac12 I_{CM}\omega_{CM}^2
\]
In assenza di forze dissipative, la conservazione dell'energia continua a valere con $K_{tot}$ al posto della sola $K_{tr}$.
\end{teorema}

\begin{nota}[Rotolamento e Attrito]
Se un oggetto ``rotola senza strisciare'' esiste implicitamente attrito. Nel \textbf{puro rotolamento}, il punto di contatto \`{e} istantaneamente fermo: l'attrito \`{e} \textbf{statico} e non compie lavoro (nessuna dissipazione di energia). Se invece c'\`{e} anche strisciamento (roto-traslazione), l'attrito \`{e} \textbf{dinamico} e dissipa parte dell'energia.
\end{nota}

\begin{teorema}[Lavoro e Potenza del Momento Torcente]
Per analogia con il caso lineare ($\Delta l = r\Delta\theta$):
\[
W = \vec\tau\cdot\Delta\vec\theta \qquad\qquad P = \frac{dW}{dt} = \vec\tau\cdot\vec\omega
\]
\end{teorema}

\begin{definizione}[Momento Angolare]
L'analogo rotazionale della quantit\`{a} di moto:
\[
\vec L = I\vec\omega
\]
\textbf{II legge di Newton generalizzata per le rotazioni} (valida anche se $I$ cambia nel tempo):
\[
\sum\vec\tau = \frac{d\vec L}{dt}
\]
che si riduce a $\sum\vec\tau=I\vec\alpha$ quando $I$ \`{e} costante.
\end{definizione}

\begin{teorema}[Conservazione del Momento Angolare]
Se $\sum\vec\tau=0$, allora $\vec L=\text{costante}$. \textbf{Attenzione}: la risultante delle forze $\sum\vec F$ pu\`{o} essere $\neq 0$, purch\'{e} il momento torcente netto sia nullo (es. forze applicate lungo l'asse di rotazione, con braccio nullo).
\end{teorema}

\begin{esempio}[Conservazione del Momento Angolare -- Pattinatrice e Tuffatore]
\textbf{Pattinatrice}: sistema isolato, $\tau=0\Rightarrow L=I\omega=\text{costante}$. Con le braccia aperte $I$ \`{e} grande e $\omega$ piccolo; raccogliendo le braccia $I$ diminuisce e $\omega$ aumenta proporzionalmente.

\textbf{Tuffatore}: non isolato (c'\`{e} la gravit\`{a}), ma il momento torcente della gravit\`{a} rispetto al CM \`{e} nullo (braccio nullo) $\Rightarrow L$ si conserva comunque durante il volo. Raccogliendo il corpo ($I$ piccolo) aumenta $\omega$ per molte rotazioni veloci; distendendosi ($I$ grande) l'atleta rallenta la rotazione per un ingresso in acqua controllato.
\end{esempio}

\begin{esempio}[Momento Angolare Vettoriale -- Persona su Piattaforma Rotante]
Se una persona ferma su una piattaforma girevole (inizialmente in quiete, $L_{tot}=0$) inizia a correre in una direzione, la piattaforma ruota nel verso \textbf{opposto}: il momento angolare della persona $L_{persona}$ deve essere bilanciato da un momento angolare uguale e opposto della piattaforma $L_{disco}$, cosicch\'{e} $L_{tot}$ resti nullo.
\end{esempio}

%=============================================================================
\section{Equilibrio in Presenza di Rotazioni e Deformazioni}
%=============================================================================

\subsection{Equilibrio Statico di un Corpo Rigido}

\begin{definizione}[Condizioni di Equilibrio]
Un corpo esteso \`{e} in equilibrio (statico) se sono verificate \textbf{entrambe}:
\[
\vec R = \sum\vec F = \vec 0 \quad \text{(traslazione)} \qquad\qquad \vec T = \sum\vec\tau = \vec 0 \quad \text{(rotazione)}
\]
6 equazioni scalari in totale (3+3). \textbf{Un corpo esteso con $\sum\vec F=0$ pu\`{o} comunque ruotare} se non anche $\sum\vec\tau=0$.
\end{definizione}

\begin{esempio}[Il Lampadario -- Tensioni Contro-Intuitive]
Un lampadario di 200 kg \`{e} sostenuto da un cavo a $60^{\circ}$ dall'orizzontale e uno orizzontale verso il muro. Dalle condizioni di equilibrio: $F_1=mg/\sin60^{\circ}\approx2260$ N, $F_2=F_1\cos60^{\circ}\approx1130$ N. Il cavo a $60^{\circ}$ deve quindi sostenere l'equivalente di $\approx231$ kg -- \textbf{pi\`{u} della massa del lampadario stesso}, un risultato controintuitivo dovuto alla scomposizione vettoriale delle tensioni.
\end{esempio}

\begin{teorema}[Principio della Leva -- Archimede]
Imponendo l'equilibrio dei momenti torcenti attorno al fulcro:
\[
\frac{r}{R} = \frac{F_P}{mg}
\]
dove $R$ \`{e} il braccio della forza applicata $F_P$ e $r$ il braccio (pi\`{u} corto) del peso da sollevare: la forza necessaria \`{e} \textbf{inversamente proporzionale} alla lunghezza del braccio della leva.
\end{teorema}

\begin{nota}[Stabilit\`{a} di un Corpo Esteso]
Inclinando un corpo appoggiato al suolo (es. un frigorifero spinto lateralmente), si genera un momento tra la reazione vincolare (al bordo di appoggio) e il peso (applicato al CM):
\begin{itemize}
  \item Piccola perturbazione: il momento del peso si oppone alla rotazione (\textbf{equilibrio stabile}, il corpo ricade in posizione).
  \item Grande perturbazione (oltre un angolo critico): il momento del peso concorda con la rotazione (\textbf{equilibrio instabile}, il corpo cade).
\end{itemize}
\textbf{Criterio generale}: un corpo appoggiato resta stabile finch\'{e} la proiezione verticale del centro di massa rimane entro la base d'appoggio.
\end{nota}

\subsection{Equilibrio Applicato al Corpo Umano}

\begin{esempio}[Ragdoll Physics in VR]
La simulazione dell'equilibrio (e squilibrio) di un corpo articolato \`{e} alla base della cosiddetta \emph{ragdoll physics}: un corpo rigido per ogni ``osso'', vincolato alle articolazioni, che collassa in modo fisicamente plausibile quando perde l'equilibrio. Evoluzione storica: dalle prime simulazioni di corpi che cadono (1997) al motore \emph{Euphoria} (met\`{a} anni 2000, con reazioni muscolari simulate, es. GTA IV) fino alle applicazioni VR in tempo reale odierne (es. SuperHotVR, 2016).
\end{esempio}

\begin{definizione}[Nomenclatura Muscolo-Scheletrica]
\textbf{Articolazione} = asse di rotazione; \textbf{inserzione} = punto di applicazione della forza muscolare; \textbf{tendine} = corda in tensione. Un muscolo pu\`{o} solo \textbf{tirare} (contraendosi): non pu\`{o} compiere lavoro positivo distendendosi, solo lavoro negativo/frenante contro un'altra forza (es. la gravit\`{a}).
\end{definizione}

\begin{esempio}[Forza del Bicipite]
Un avambraccio orizzontale regge una sfera. Poich\'{e} la 1\textsuperscript{a} condizione di equilibrio introduce l'incognita aggiuntiva della reazione vincolare al gomito $F_J$, conviene usare \emph{solo} la 2\textsuperscript{a} condizione (momenti attorno al gomito, dove $F_J$ ha braccio nullo):
\[
F_M = \frac{(0.15\,m_A+0.35\,m)g}{0.05} \approx 402\ \text{N}
\]
($m_A=2.0$ kg massa avambraccio a 15 cm dal gomito, $m=5.0$ kg sfera a 35 cm, inserzione del bicipite a 5 cm). Se il braccio si piega di un angolo $\theta$, \textbf{tutti} i bracci si riducono dello stesso fattore $\cos\theta$, che si semplifica nell'equazione dei momenti: $F_M$ resta \textbf{invariata}. Poich\'{e} $F_M \propto 1/d_{inserzione}$, una variazione di soli 5 mm nel punto di inserzione del tendine cambia la forza muscolare richiesta di circa il 10\%.
\end{esempio}

\begin{esempio}[Il Carico sulla Colonna Vertebrale]
Per una persona china in avanti (busto a $60^{\circ}$ dall'orizzontale), imponendo l'equilibrio dei momenti attorno alla vertebra L5 (fulcro del busto) con i pesi di testa, braccia e tronco (rispettivamente $0.07w$, $0.12w$, $0.46w$ del peso corporeo totale $w$) si ottiene una tensione dei muscoli interspinali $F_M\approx2.2w$, e una forza risultante sulla vertebra $F_V\approx2.5w$. \textbf{Sollevare un peso senza inclinare la schiena} (mantenendo il busto verticale) evita di moltiplicare per 2.5 (o pi\`{u}, aggiungendo il peso sollevato) il carico sulla colonna vertebrale.
\end{esempio}

\subsection{Deformazioni Elastiche}

\begin{definizione}[Regime Elastico -- Legge di Hooke Generalizzata]
Una forza applicata a un corpo (non pi\`{u} rigido) ne causa una deformazione. Per piccole deformazioni vale la proporzionalit\`{a} $F=k\Delta L$ (\textbf{regime elastico}): rimuovendo la forza, il corpo torna alla forma originale.
\end{definizione}

\begin{center}
\begin{tikzpicture}
\begin{axis}[width=9.5cm, height=6cm, xlabel={Allungamento $\Delta L$}, ylabel={Forza $F$}, xtick=\empty, ytick=\empty, axis lines=left, font=\small, xmin=0, ymin=0, xmax=4.3, ymax=6.5]
\addplot[domain=0:2, samples=50, thick, blue] {2.2*x};
\addplot[domain=2:3.6, samples=50, thick, blue] {4.4+1.1*(x-2)-0.15*(x-2)^2};
\addplot[mark=*, mark size=2pt, blue] coordinates {(3.6,5.6)};
\draw[dashed] (axis cs:2,0) -- (axis cs:2,4.4);
\node[align=center] at (axis cs:1,1.2) {\footnotesize regione\\\footnotesize elastica};
\node[align=center] at (axis cs:2.8,1.6) {\footnotesize regione\\\footnotesize plastica};
\node[above] at (axis cs:2,4.4) {\footnotesize limite elastico};
\node[anchor=south west] at (axis cs:3.6,5.6) {\footnotesize rottura};
\end{axis}
\end{tikzpicture}
\captionof{figure}{Curva forza-allungamento: nella regione elastica (proporzionale, legge di Hooke) la deformazione \`{e} reversibile; oltre il limite elastico si entra nella regione plastica (deformazione permanente), fino al punto di rottura.}
\end{center}

\begin{teorema}[Modulo di Young, Sforzo e Deformazione -- Tensione/Compressione]
Per un cilindro di lunghezza $L_0$ e sezione $A$, sottoposto a trazione (o compressione) longitudinale:
\[
\Delta L = \frac{1}{E}\frac{F}{A}L_0 \qquad\qquad E = \frac{\text{sforzo}}{\text{deformazione}} = \frac{F/A}{\Delta L/L_0}
\]
dove $F/A$ (N/m$^2$) \`{e} lo \textbf{sforzo} (\emph{stress}) e $\Delta L/L_0$ (numero puro) \`{e} la \textbf{deformazione} (\emph{strain}). Un materiale rigido ha $E$ grande; uno duttile ha $E$ piccolo.
\end{teorema}

\begin{teorema}[Sforzo di Taglio]
Per forze applicate \textbf{tangenzialmente} alle facce di un solido (anzich\'{e} perpendicolarmente):
\[
\Delta L = \frac{1}{G}\frac{F}{A}L_0
\]
dove $G$ (\textbf{modulo di taglio}) gioca il ruolo di $E$, e $A$ \`{e} ora l'area della faccia \emph{parallela} alla forza applicata (non quella perpendicolare come per tensione/compressione). Regola pratica: $G$ \`{e} generalmente il 30-50\% di $E$.
\end{teorema}

\begin{teorema}[Modulo di Compressione Volumetrica -- Bulk Modulus]
Per una pressione applicata su tutto il volume di un corpo (es. immerso in un fluido):
\[
\frac{\Delta V}{V_0} = -\frac{1}{B}\Delta P
\]
Il segno negativo indica che un aumento di pressione \emph{riduce} il volume.
\end{teorema}

\begin{center}
\small
\begin{tabular}{lccc}
\toprule
\textbf{Materiale} & $E$ (10$^9$ N/m$^2$) & $G$ (10$^9$ N/m$^2$) & $B$ (10$^9$ N/m$^2$) \\
\midrule
Acciaio & 200 & 80 & 140 \\
Alluminio & 70 & 25 & 70 \\
Osso (arto) & 15 & 80 & -- \\
Legno (pino), parallelo alle fibre & 10 & -- & -- \\
Legno (pino), perpendicolare alle fibre & 1 & -- & -- \\
Cemento & 20 & trascurabile & -- \\
Acqua & -- & -- & 2.0 \\
\bottomrule
\end{tabular}
\captionof{table}{Moduli elastici tipici (valori indicativi). Il legno \`{e} molto pi\`{u} rigido lungo le fibre che trasversalmente; cemento/laterizio hanno modulo di taglio trascurabile (materiali fragili al taglio); l'osso ha un modulo di taglio sorprendentemente alto rispetto a $E$.}
\end{center}

\begin{nota}[Massimo Sforzo, Frattura e Fattore di Sicurezza]
Superata una soglia di sforzo massimo, il materiale si frattura (per tensione, compressione o taglio); \`{e} difficile prevedere con precisione dove e quando. Gli \textbf{sforzi di tensione e di taglio sono generalmente i pi\`{u} pericolosi} (molti materiali fragili, come cemento e pietra, non tollerano quasi nessuno sforzo di trazione o taglio). Per questo motivo i progetti architettonici/civili impongono per legge un \textbf{fattore di sicurezza} (rapporto tra sforzo massimo tollerabile e sforzo effettivo) tipicamente tra \textbf{3 e 10}.
\end{nota}

%=============================================================================
\section{Onde Meccaniche}
%=============================================================================

\subsection{Concetti Fondamentali}

\begin{definizione}[Onda Meccanica]
La propagazione, in un mezzo, di una perturbazione dell'equilibrio: una successione infinita di moti armonici (cfr. \S5) che si trasmette da un punto all'altro del mezzo grazie a forze di coesione interne.
\end{definizione}

\begin{teorema}[Cosa Trasporta un'Onda]
In un'onda meccanica \textbf{ideale} non si ha mai trasporto di \textbf{materia}, ma solo di \textbf{energia}: ogni punto del mezzo oscilla attorno alla propria posizione di equilibrio senza spostarsi nel verso di propagazione (es. una boa sull'acqua sale e scende al passaggio di un'onda, ma non viene trascinata orizzontalmente).
\end{teorema}

\begin{definizione}[Grandezze di un'Onda Periodica]
\textbf{Ampiezza} $A$: elongazione massima. \textbf{Lunghezza d'onda} $\lambda$: distanza spaziale tra due punti identici (stessa fase). \textbf{Periodo} $T$, \textbf{frequenza} $f=1/T$.

Un'onda ha \textbf{due velocit\`{a} distinte}: la velocit\`{a} \textbf{materiale} (oscillazione di ciascun punto, $v_P=\partial y/\partial t$) e la velocit\`{a} \textbf{d'onda} (propagazione dell'energia):
\[
v = \lambda f = \frac{\lambda}{T}
\]
\end{definizione}

\begin{definizione}[Onde Trasversali e Longitudinali]
\textbf{Trasversali}: lo spostamento della materia \`{e} \textbf{perpendicolare} alla propagazione (es. corda scossa verticalmente). \textbf{Longitudinali}: lo spostamento \`{e} \textbf{parallelo} alla propagazione, con zone alternate di compressione ed espansione (es. il suono).
\end{definizione}

\begin{teorema}[Velocit\`{a} di Propagazione]
In ogni caso, velocit\`{a} = (fattore elastico)/(fattore inerziale):
\[
v_{\text{corda}} = \sqrt{\frac{F_T}{\mu}} \qquad v_{\text{solidi}} = \sqrt{\frac{E}{\rho}} \qquad v_{\text{liquidi/gas}} = \sqrt{\frac{B}{\rho}}
\]
dove $F_T$ = tensione, $\mu=M/L$ = densit\`{a} lineare, $E$ = modulo di Young, $B$ = modulo di compressione (cfr. \S8).
\end{teorema}

\begin{nota}[Perch\'{e} i liquidi/gas non supportano onde trasversali]
Solo i \textbf{solidi} resistono a sforzi di taglio (modulo $G$, cfr. \S8) e possono quindi propagare sia onde longitudinali che trasversali. I \textbf{liquidi e i gas} non hanno forma propria e non resistono al taglio: supportano solo onde \textbf{longitudinali} (di compressione).
\end{nota}

\begin{esempio}[Onde Sismiche e lo Stato del Nucleo Terrestre]
Un terremoto genera sia onde \textbf{P} (\emph{pressure}, longitudinali) che onde \textbf{S} (\emph{shear}, trasversali). Le stazioni sismiche rivelano una \textbf{zona d'ombra} per le onde S (assenti tra $103^{\circ}$ e $180^{\circ}$ dall'epicentro, misurati lungo la superficie terrestre): poich\'{e} le onde S non attraversano un mezzo liquido, questo dimostra che il \textbf{nucleo esterno terrestre \`{e} allo stato liquido}.
\end{esempio}

\subsection{Energia, Funzione d'Onda, Sovrapposizione}

\begin{teorema}[Energia e Intensit\`{a} di un'Onda]
L'energia trasportata da un'onda \`{e} proporzionale al \textbf{quadrato dell'ampiezza} (conseguenza dell'energia dell'oscillatore armonico, $E\propto A^2$, cfr. \S5). L'\textbf{intensit\`{a}} \`{e} potenza per unit\`{a} di area (perpendicolare alla propagazione). Per un'onda sferica, l'energia si distribuisce su una superficie $4\pi r^2$:
\[
I = \frac{P}{4\pi r^2} \quad\Rightarrow\quad \frac{I_1}{I_2}=\frac{r_2^2}{r_1^2} \quad\Rightarrow\quad \frac{A_1}{A_2}=\frac{r_2}{r_1}
\]
(conservazione dell'energia: l'intensit\`{a} decresce come $1/r^2$, l'ampiezza come $1/r$).
\end{teorema}

\begin{definizione}[Funzione d'Onda 1D]
\[
u(x,t) = A\cos(kx\pm\omega t+\varphi_0) \qquad k=\frac{2\pi}{\lambda}\ \text{(numero d'onda)} \qquad \omega=\frac{2\pi}{T}
\]
\textbf{Attenzione}: il numero d'onda $k$ non ha nulla a che vedere con la costante elastica $k$ della molla (\S5) -- stesso simbolo, grandezza fisica completamente diversa.
\end{definizione}

\begin{teorema}[Principio di Sovrapposizione]
Due o pi\`{u} onde meccaniche possono passare per lo stesso punto agendo indipendentemente l'una dall'altra; lo spostamento risultante \`{e} la \textbf{somma} degli spostamenti individuali:
\[
u_{tot} = u_1+u_2+\ldots
\]
Vale quando la relazione tra deformazione e forza di richiamo \`{e} \textbf{lineare} (regime elastico/armonico); smette di valere per perturbazioni troppo grandi (oltre il limite elastico, cfr. \S8).
\end{teorema}

\subsection{Riflessione, Trasmissione, Interferenza}

\begin{teorema}[Riflessione agli Estremi]
Per un'onda su una corda che incontra un'interfaccia (es. un'estremit\`{a} fissa o libera):
\begin{itemize}
  \item \textbf{Estremo fisso} ($u=0$ sempre): riflessione \textbf{con sfasamento} ($A_R=-A_I$, l'onda riflessa \`{e} invertita).
  \item \textbf{Estremo libero} ($\partial u/\partial x=0$): riflessione \textbf{senza sfasamento} ($A_R=A_I$).
\end{itemize}
\end{teorema}

\begin{teorema}[Trasmissione tra Due Mezzi]
All'interfaccia tra due mezzi con numero d'onda $k_I$ (incidente) e $k_T$ (trasmesso), imponendo continuit\`{a} di spostamento e di pendenza (tensione bilanciata):
\[
A_T = A_I\,\frac{2k_I}{k_I+k_T} \qquad\qquad A_R = A_I\,\frac{k_I-k_T}{k_I+k_T}
\]
\end{teorema}

\begin{definizione}[Interferenza e Battimenti]
Due onde \textbf{coerenti} (stessa $f$ e $\lambda$, sfasamento costante) generano \textbf{interferenza}: costruttiva se in fase ($\Delta\varphi\approx0$), distruttiva se in opposizione di fase ($\Delta\varphi\approx\pi$). Applicazioni: cuffie anti-rumore, progettazione acustica di sale. Se le frequenze sono leggermente diverse ($f_1\approx f_2$), si generano \textbf{battimenti}: un'onda a frequenza intermedia, la cui ampiezza \`{e} modulata a una frequenza $|f_1-f_2|$ (usato per accordare strumenti a corda per confronto con un diapason).
\end{definizione}

\subsection{Onde Stazionarie}

\begin{teorema}[Onde Stazionarie]
La sovrapposizione di un'onda progressiva e una regressiva identiche genera un'\textbf{onda stazionaria}:
\[
u_I+u_R = 2A\cos(kx)\cos(\omega t)
\]
Lo spazio e il tempo si separano matematicamente: \textbf{un'onda stazionaria non \`{e} un'onda} (non trasporta n\'{e} materia n\'{e} energia nel senso di propagazione), ma il prodotto dell'interferenza tra due onde reali. I \textbf{nodi} hanno sempre ampiezza nulla; i \textbf{ventri} oscillano tra $0$ e $2A$. L'energia resta immagazzinata nel mezzo, oscillando tra energia cinetica (massima quando il mezzo passa per l'equilibrio) e potenziale elastica (massima alla massima elongazione).
\end{teorema}

\begin{teorema}[Armoniche di una Corda con Estremi Fissi]
Solo specifiche lunghezze d'onda producono onde stazionarie stabili su una corda di lunghezza $L$ fissata a entrambi gli estremi (che devono essere \textbf{nodi}):
\[
\lambda_n = \frac{2L}{n} \qquad f_n = \frac{nv}{2L} \qquad (n=1,2,3,\ldots)
\]
$n=1$: \textbf{fondamentale} ($\lambda=2L$); $n>1$: \textbf{armoniche superiori}.
\end{teorema}

\begin{center}
\begin{tikzpicture}[font=\small]
\foreach \n/\y in {1/0,2/-1.6,3/-3.2} {
  \draw[thick] (0,\y) -- (0,\y+0.35) (0,\y-0.35) -- (0,\y);
  \draw[thick] (5,\y) -- (5,\y+0.35) (5,\y-0.35) -- (5,\y);
  \draw[thick, blue, domain=0:5, samples=100, smooth] plot ({\x}, {\y+0.5*sin(\n*180*\x/5)});
  \draw[thick, blue, dashed, domain=0:5, samples=100, smooth] plot ({\x}, {\y-0.5*sin(\n*180*\x/5)});
  \node[left] at (-0.3,\y) {$n=\n$};
}
\end{tikzpicture}
\captionof{figure}{Le prime tre armoniche di una corda fissata a entrambi gli estremi: $n=1$ (fondamentale, un ventre), $n=2$ (due semi-lunghezze d'onda, un nodo centrale), $n=3$ (tre semi-lunghezze d'onda). Le curve continua e tratteggiata rappresentano la corda a due istanti diversi durante l'oscillazione.}
\end{center}

%=============================================================================
\section{Acustica e Percezione Uditiva}
%=============================================================================

\subsection{Il Suono come Onda Longitudinale}

\begin{definizione}[Suono]
Onda meccanica \textbf{longitudinale} di pressione. \textbf{Senza un mezzo materiale il suono non pu\`{o} propagarsi} -- un errore fisico comune in film e videogiochi ambientati nello spazio, dove si sentono esplosioni nel vuoto.
\end{definizione}

\begin{center}
\small
\begin{tabular}{lc}
\toprule
\textbf{Mezzo (a 20	extdegree{}C, 1 atm)} & \textbf{Velocit\`{a} del suono (m/s)} \\
\midrule
Aria (0	extdegree{}C) & 331 \\
Aria (20	extdegree{}C) & 343 \\
Elio & 1005 \\
Acqua & 1440 \\
Ferro e acciaio & $\approx5000$ \\
\bottomrule
\end{tabular}
\captionof{table}{Velocit\`{a} del suono in diversi mezzi. In aria: $v=(331+0.6T)$ m/s, $T$ in 	extdegree{}C.}
\end{center}

\begin{nota}[Spostamento e pressione sono sfasati di $90^{\circ}$]
Nell'onda sonora, lo spostamento delle molecole \`{e} sfasato di $\pi/2$ rispetto all'onda di pressione: a un massimo (o minimo) di pressione corrisponde spostamento nullo, e viceversa. Le due rappresentazioni sono equivalenti e intercambiabili a seconda del problema.
\end{nota}

\begin{definizione}[Tono e Intervallo Udibile]
Il \textbf{tono} (\emph{pitch}) \`{e} l'equivalente percettivo della frequenza. L'orecchio umano \`{e} sensibile tra circa \textbf{20 Hz e 20 000 Hz}; sotto i 20 Hz si parla di \textbf{infrasuoni}, sopra i 20 kHz di \textbf{ultrasuoni} (es. cane $\sim$50 kHz, pipistrello $\sim$100 kHz, usato per l'ecolocalizzazione).
\end{definizione}

\begin{teorema}[Teorema di Nyquist e Campionamento Digitale]
Un campionamento a frequenza $f_s$ registra/riproduce correttamente onde fino a $f_s/2$. Frequenze superiori vengono erroneamente interpretate come frequenze pi\`{u} basse (\textbf{aliasing}); la soluzione \`{e} filtrare il segnale (filtro passa-basso) prima del campionamento. L'audio digitale standard campiona a \textbf{44.1 kHz} ($f_s/2=22.05$ kHz, appena sopra la sensibilit\`{a} umana).
\end{teorema}

\subsection{Intensit\`{a}, Decibel, Percezione}

\begin{definizione}[Scala dei Decibel]
Data l'ampiezza dell'intervallo di udibilit\`{a} (12 ordini di grandezza in intensit\`{a}), si usa una scala logaritmica:
\[
\beta\ (\text{dB}) = 10\log\frac{I}{I_0} \qquad I_0 = 10^{-12}\ \text{W/m}^2 \ \text{(soglia dell'udito)}
\]
\end{definizione}

\begin{center}
\small
\begin{tabular}{lc}
\toprule
\textbf{Sorgente} & \textbf{Livello (dB)} \\
\midrule
Aviogetto a 30 m / soglia del dolore & 120--140 \\
Concerto rock al chiuso & 120 \\
Strada trafficata & 70 \\
Conversazione ordinaria (50 cm) & 65 \\
Sussurro & 20 \\
Soglia dell'udito & 0 \\
\bottomrule
\end{tabular}
\captionof{table}{Livelli di intensit\`{a} sonora tipici. La maggior parte dei fenomeni udibili copre circa 100 dB.}
\end{center}

\begin{nota}[Rumore di fondo e rischio uditivo]
In condizioni tipiche odierne (traffico, guida) il rumore di fondo \`{e} gi\`{a} 75-80 dB; per udire chiaramente una conversazione o musica in cuffia in quelle condizioni serve aggiungere altri $\sim$40 dB, portando il livello vicino alla \textbf{soglia del dolore} (120 dB) -- da cui il rischio di danno uditivo (spesso irreversibile) in ambienti gi\`{a} rumorosi.
\end{nota}

\begin{teorema}[Curve Isofoniche -- Risposta in Frequenza dell'Orecchio]
La sensibilit\`{a} dell'orecchio umano \textbf{non \`{e} uniforme} sulle frequenze udibili: la soglia dell'udito ha un minimo (massima sensibilit\`{a}) attorno a \textbf{2-5 kHz} (frequenze tipiche delle armoniche del parlato), e cresce ripidamente sotto i 100 Hz e sopra i 10 kHz. Ad esempio, $62$ dB a $100$ Hz produce la stessa sensazione uditiva di $40$ dB a $1$ kHz. Questo richiede una correzione non lineare (\textbf{loudness control}) nei sistemi audio.
\end{teorema}

\subsection{Musica: Armoniche, Timbro, Spettro}

\begin{definizione}[Note e Armoniche di una Corda]
Su una corda vibrante fissata a entrambi gli estremi (\S9), la fondamentale $f_1$ definisce la nota; le armoniche superiori sono multipli interi: $f_n=nf_1$. Il \textbf{tono} (nota percepita) \`{e} dato dalla fondamentale; il \textbf{timbro} \`{e} dato dalla distribuzione di ampiezza tra le armoniche superiori, ed \`{e} ci\`{o} che distingue il suono di strumenti diversi che suonano la stessa nota.
\end{definizione}

\begin{teorema}[Teorema di Fourier e Principio di Sovrapposizione in Musica]
Qualsiasi funzione periodica (quindi qualsiasi suono complesso) pu\`{o} essere scritta come somma di componenti sinusoidali (armoniche) a frequenze multiple della fondamentale, ciascuna con una propria ampiezza. Lo \textbf{spettro} di uno strumento (intensit\`{a} delle armoniche in funzione della frequenza) \`{e} l'analisi matematica del suo timbro.
\end{teorema}

\begin{definizione}[Strumenti a Fiato -- Tubo Aperto vs. Chiuso]
Ragionando in termini di \textbf{pressione}, agli estremi aperti la pressione non pu\`{o} variare rispetto all'ambiente esterno: devono essere \textbf{nodi} di pressione.
\begin{itemize}
  \item \textbf{Tubo aperto-aperto}: $L=n\frac{\lambda_n}{2}$, quindi $f_n = \dfrac{nv}{2L}$ -- supporta tutte le armoniche.
  \item \textbf{Tubo aperto-chiuso}: l'estremo chiuso \`{e} un \textbf{ventre} di pressione; $L=(2n-1)\frac{\lambda_n}{4}$, quindi $f_n = \dfrac{(2n-1)v}{4L}$ -- supporta \textbf{solo le armoniche dispari}.
\end{itemize}
\end{definizione}

\begin{definizione}[Rumore e i suoi ``Colori'']
Combinazione di frequenze in rapporti casuali (nessun tono dominante udibile). \textbf{Rumore bianco}: spettro piatto (tutte le frequenze con la stessa ampiezza). \textbf{Rumore rosa}: ampiezza $\propto 1/f$. \textbf{Rumore grigio}: rumore bianco corretto per le curve isofoniche (percezione costante su tutte le frequenze).
\end{definizione}

\begin{teorema}[Interferenza e Battimenti in Acustica]
Onde sonore coerenti interferiscono costruttivamente (in fase) o distruttivamente (in opposizione di fase) -- principio alla base delle \textbf{cuffie anti-rumore} e della progettazione acustica di sale da concerto (es. il canale centrale di un sistema 5.1, dedicato al parlato, va posizionato per evitare interferenza distruttiva). Due frequenze vicine ($f_1\approx f_2$) generano \textbf{battimenti} a frequenza $|f_1-f_2|$, usati per accordare strumenti a corda per confronto (es. con un diapason).
\end{teorema}

\begin{nota}[Risoluzione Temporale dell'Orecchio Umano]
La capacit\`{a} di distinguere due impulsi sonori ravvicinati dipende dalla durata dell'impulso: per impulsi lunghi ($\sim$80 ms) la risoluzione \`{e} migliore ($\sim$3 ms), per impulsi brevi ($\sim$10 ms) \`{e} peggiore ($\sim$20-30 ms). Conseguenza pratica per la VR: la generazione del suono a livello di un singolo frame (a 75 fps, $\sim$13 ms) \`{e} coerente con questa risoluzione temporale; inoltre, un battimento con $|f_1-f_2|\geq30$ Hz \`{e} difficilmente percepito come tale.
\end{nota}

\subsection{Effetto Doppler}

\begin{teorema}[Effetto Doppler]
Quando sorgente e/o osservatore sono in moto relativo lungo la congiungente, la frequenza percepita $f_o$ differisce dalla frequenza emessa $f_s$:
\[
f_o = f_s\,\frac{v\pm v_o}{v\mp v_s}
\]
dove $v$ \`{e} la velocit\`{a} del suono nel mezzo, $v_o,v_s$ le velocit\`{a} di osservatore e sorgente (con segno secondo la convenzione: positive se avvicinano sorgente e osservatore). \textbf{L'effetto non \`{e} simmetrico}: a parit\`{a} di velocit\`{a} relativa, muovere la sorgente o muovere l'osservatore produce risultati quantitativamente diversi (sebbene qualitativamente analoghi).
\end{teorema}

\begin{esempio}[Il ``Muro del Suono'' e la Sirena dell'Ambulanza]
Se la velocit\`{a} della sorgente eguaglia quella del suono, i fronti d'onda emessi si sovrappongono tutti davanti alla sorgente in un'unica interferenza costruttiva di altissima energia: il \textbf{muro del suono}. Nel caso pi\`{u} comune (auto/ambulanza che ci sorpassa), nell'istante in cui la sorgente ci supera il segno di $v_s$ cambia bruscamente, causando la caratteristica \textbf{discontinuit\`{a}} nel tono percepito della sirena.
\end{esempio}

\begin{esempio}[Sensori di Distanza in VR: Ultrasuoni vs. Infrarossi]
Un sensore a ultrasuoni impiega $\sim$6 ms per percorrere 1 m (andata, quindi 12 ms andata/ritorno). Con un target di \textbf{75 fps} in VR, un frame dura $1/75\approx13.3$ ms: la sola misura ultrasonica di andata/ritorno consumerebbe quasi l'intero budget di tempo per frame, lasciando poco margine per l'elaborazione -- e la velocit\`{a} del suono dipende inoltre dalla temperatura, introducendo un ulteriore errore sistematico se non corretta. Per questo i visori VR utilizzano sensori a \textbf{infrarossi} (onde elettromagnetiche, velocit\`{a} $\sim10^6$ volte quella del suono): il tempo di volo diventa trascurabile rispetto al budget del frame, al costo di richiedere una risoluzione temporale $10^6$ volte superiore (nanosecondi anzich\'{e} millisecondi) per misurare comunque la distanza con precisione.
\end{esempio}

%=============================================================================
\section{Ottica Geometrica, Strumenti Ottici e Visori VR}
%=============================================================================

\subsection{Ottica Geometrica: Concetti di Base}

\begin{definizione}[Ottica Geometrica vs. Ondulatoria]
La luce \`{e} un'onda elettromagnetica (non necessita di un mezzo per propagarsi). L'\textbf{ottica geometrica} considera solo la \emph{direzione} di propagazione (\textbf{raggi}, linee perpendicolari ai fronti d'onda) ed \`{e} alla base del rendering in tempo reale; l'\textbf{ottica ondulatoria} descrive fenomeni di fase (interferenza, diffrazione), spesso troppo costosi da calcolare in tempo reale e quindi approssimati o pre-calcolati offline.
\end{definizione}

\begin{definizione}[Riflessione Speculare e Diffusiva]
\textbf{Speculare}: superficie piana, stessa normale ovunque, angolo di incidenza = angolo di riflessione ($\theta_i=\theta_r$) rispetto alla normale. \textbf{Diffusiva}: superficie scabra, la normale cambia punto per punto, ma la legge di riflessione \`{e} comunque rispettata localmente in ogni punto. \textbf{La riflessione diffusiva \`{e} alla base della visione}: solo le superfici diffusive sono visibili da qualunque angolazione (Galileo us\`{o} questo principio per dedurre che la Luna, riflettendo diffusamente e non specularmente, non poteva essere una sfera perfetta).
\end{definizione}

\subsection{Specchi}

\begin{teorema}[Specchi Piani]
Per congruenza di triangoli, l'immagine virtuale di uno specchio piano si forma a distanza $d_i=d_o$ dietro lo specchio, con le stesse dimensioni dell'oggetto (nessun ingrandimento). Uno specchio piano \textbf{non ribalta} realmente sinistra/destra: la percezione di inversione laterale \`{e} un effetto della nostra esperienza percettiva, non della fisica dello specchio.
\end{teorema}

\begin{definizione}[Specchi Sferici]
\textbf{Concavo} (curvatura verso l'interno): converge i raggi paralleli in un fuoco. \textbf{Convesso} (curvatura verso l'esterno): diverge i raggi; il fuoco \`{e} virtuale (dietro lo specchio). Nell'approssimazione \textbf{parassiale} (raggi vicini all'asse, piccoli angoli -- detta anche ottica gaussiana), lo specchio converge (quasi) tutti i raggi paralleli in un unico punto, evitando l'\textbf{aberrazione sferica} (la naturale tendenza di uno specchio sferico a formare una zona di confusione anzich\'{e} un punto, per raggi lontani dall'asse).
\end{definizione}

\begin{teorema}[Equazione degli Specchi (di Gauss)]
Per similitudine di triangoli (raggio parallelo all'asse che riflette per il fuoco; raggio che passa per il fuoco e riflette parallelo all'asse):
\[
\frac{1}{f} = \frac{1}{d_i}+\frac{1}{d_o} \qquad f = \frac{r}{2} \qquad m = \frac{h_i}{h_o} = -\frac{d_i}{d_o}
\]
Valida in generale (non solo per oggetti tra $C$ e $F$), con la \textbf{convenzione dei segni}: $r,f>0$ per specchi concavi ($r,f<0$ per convessi); $d_i>0$ se l'immagine \`{e} reale (davanti allo specchio); $h_i<0$ se l'immagine \`{e} capovolta.
\end{teorema}

\begin{center}
\begin{tikzpicture}[font=\small, scale=1]
\draw[thick] plot[domain=-1.6:1.6, samples=50] ({-0.15*\x*\x},{\x});
\draw[thick, dashed] (-3.6,0) -- (0.4,0);
\node[below] at (-2.6,-0.05) {$C$};
\node[below] at (-1.3,-0.05) {$F$};
\node[below] at (0,-0.05) {$A$};
\draw[thick] (-3.3,0) -- (-3.3,1.1);
\draw[-Stealth, thick, red] (-3.3,1.1) -- (0,1.1) -- (-2.145,-0.715);
\draw[-Stealth, thick, blue] (-3.3,1.1) -- (0,-0.715) -- (-2.145,-0.715);
\node[above] at (-3.3,1.1) {O};
\draw[thick, dashed] (-2.145,0) -- (-2.145,-0.715);
\node[below] at (-2.145,-0.75) {I};
\end{tikzpicture}
\captionof{figure}{Specchio concavo: il raggio parallelo all'asse (rosso) riflette passando per il fuoco $F$; il raggio che passa per $F$ (blu) riflette parallelo all'asse. L'intersezione dei due raggi riflessi individua l'immagine $I$, reale e capovolta.}
\end{center}

\subsection{Rifrazione}

\begin{definizione}[Indice di Rifrazione]
\[
n = \frac{c}{v} \qquad c\approx3\times10^8\ \text{m/s (vuoto)}
\]
\end{definizione}

\begin{center}
\small
\begin{tabular}{lc}
\toprule
\textbf{Mezzo} & $n$ (a $\lambda=589$ nm) \\
\midrule
Vuoto & 1.0000 \\
Aria & 1.0003 \\
Acqua & 1.33 \\
Vetro crown & 1.52 \\
Diamante & 2.42 \\
\bottomrule
\end{tabular}
\end{center}

\begin{teorema}[Legge di Snell]
\[
n_1\sin\theta_1 = n_2\sin\theta_2
\]
Il raggio si avvicina alla normale entrando in un mezzo a velocit\`{a} di propagazione minore (n maggiore), se ne allontana nel caso opposto.
\end{teorema}

\begin{nota}[Dispersione e Prisma]
L'indice di rifrazione dipende (debolmente) dalla lunghezza d'onda ($n=n(\lambda)$, tipicamente decrescente con $\lambda$): un prisma separa spazialmente la luce bianca nello spettro visibile, con il \textbf{rosso deviato meno} e il \textbf{violetto deviato di pi\`{u}}.
\end{nota}

\begin{teorema}[Riflessione Totale Interna]
Passando da un mezzo pi\`{u} denso otticamente a uno meno denso ($n_1>n_2$), oltre un \textbf{angolo critico} $\theta_c$ tutta la luce viene riflessa internamente (nessuna trasmissione):
\[
\sin\theta_c = \frac{n_2}{n_1}
\]
(es. acqua-aria: $\theta_c\approx49^{\circ}$; vetro-aria: $\theta_c\approx36$-$39^{\circ}$). Applicazioni: prismi nei binocoli (riflessione al 100\%, a differenza di uno specchio metallico); \textbf{fibre ottiche} (core ad alto $n$, cladding a $n$ minore, propagazione per riflessione totale anche in curve strette).
\end{teorema}

\subsection{Lenti Sottili}

\begin{definizione}[Lenti Convergenti e Divergenti]
Una lente (di vetro/plastica) rifrange la luce due volte (entrata e uscita). \textbf{Convergente} (pi\`{u} spessa al centro): converge raggi paralleli in un fuoco reale oltre la lente. \textbf{Divergente} (pi\`{u} spessa ai bordi): diverge i raggi; il fuoco \`{e} virtuale. Il \textbf{piano focale} \`{e} il piano perpendicolare all'asse passante per $F$: vi convergono tutti i raggi paralleli provenienti dall'infinito, non solo quelli paralleli all'asse ottico.
\end{definizione}

\begin{teorema}[Equazione delle Lenti Sottili]
Stessa forma dell'equazione degli specchi, con le distanze misurate dal \textbf{centro} della lente:
\[
\frac{1}{f} = \frac{1}{d_i}+\frac{1}{d_o} \qquad m = \frac{h_i}{h_o} = -\frac{d_i}{d_o} \qquad P = \frac{1}{f}\ \text{(potenza, in diottrie D, $f$ in metri)}
\]
$P>0$ per lenti convergenti, $P<0$ per lenti divergenti. Per sistemi di pi\`{u} lenti, l'immagine della prima diventa l'oggetto della seconda; l'ingrandimento totale \`{e} il prodotto degli ingrandimenti singoli.
\end{teorema}

\begin{center}
\begin{tikzpicture}[font=\small]
\draw[thick, blue!40!black] (0,-1.8) to[out=80,in=-80] (0,1.8) to[out=100,in=260] (-0.25,1.8) to[out=-80,in=100] (-0.25,-1.8) to[out=80,in=-100] cycle;
\draw[thick, dashed] (-3.5,0) -- (3.5,0);
\node[below] at (-1.2,-0.05) {$F'$};
\node[below] at (1.2,-0.05) {$F$};
\node[below] at (0,-0.05) {$A$};
\draw[thick] (-2.5,0) -- (-2.5,1.0);
\node[above] at (-2.5,1.0) {O};
\draw[-Stealth, thick, red] (-2.5,1.0) -- (0,1.0) -- (2.31,-0.92);
\draw[-Stealth, thick, teal] (-2.5,1.0) -- (-1.2,0) -- (0,-0.92) -- (2.31,-0.92);
\draw[thick, dashed] (2.31,0) -- (2.31,-0.92);
\node[below] at (2.31,-0.95) {I};
\end{tikzpicture}
\captionof{figure}{Lente convergente: il raggio parallelo all'asse (rosso) rifrange passando per il fuoco $F$; il raggio che passa per il fuoco anteriore $F'$ (verde) rifrange parallelo all'asse. La loro intersezione forma l'immagine reale $I$, capovolta.}
\end{center}

\subsection{L'Occhio Umano e la Macchina Fotografica}

\begin{definizione}[Anatomia Ottica dell'Occhio]
La \textbf{cornea} ($n=1.376$) fornisce la maggior parte della rifrazione iniziale; il \textbf{cristallino}, con curvatura regolata dai muscoli ciliari (\textbf{accomodazione}), mette a fuoco sulla \textbf{retina}; l'\textbf{iride} regola, tramite la \textbf{pupilla}, la quantit\`{a} di luce entrante.
\end{definizione}

\begin{nota}[Difetti Visivi Comuni]
\textbf{Miopia}: il punto di convergenza cade \emph{davanti} alla retina (a riposo); si corregge con una lente \textbf{divergente}. \textbf{Ipermetropia}: il punto di convergenza cade \emph{dietro} la retina; si corregge con una lente \textbf{convergente}. \textbf{Astigmatismo}: distorsione (un punto diventa una linea), corretto con lenti cilindriche.
\end{nota}

\begin{esempio}[La Macchina Fotografica come ``Camera Virtuale'']
Il modello ottico di una macchina fotografica (obiettivo con apertura regolabile dal diaframma, sensore) \`{e} lo stesso usato per le \textbf{camere virtuali in computer grafica 3D}. Il \textbf{rapporto di apertura} (f-stop) \`{e} $f/D$ (focale/diametro apertura); aperture maggiori (f-stop minore) producono una \textbf{profondit\`{a} di campo} minore (regione a fuoco pi\`{u} sottile), con maggiore sfocatura (cerchio di confusione pi\`{u} grande) fuori fuoco.
\end{esempio}

\begin{definizione}[Aberrazioni delle Lenti]
Oltre all'\textbf{aberrazione sferica} (gi\`{a} vista per gli specchi), le lenti soffrono di: \textbf{aberrazione cromatica} (indici di rifrazione diversi per $\lambda$ diverse $\Rightarrow$ fuochi diversi per colori diversi, corretta con un \textbf{doppietto acromatico} di due vetri/forme diverse); \textbf{coma}, \textbf{astigmatismo} e \textbf{curvatura di campo} (per punti fuori asse).
\end{definizione}

\subsection{Ottica Applicata ai Visori VR}

\begin{teorema}[Risoluzione Angolare]
La capacit\`{a} di distinguere due punti separati da un angolo $\theta$. L'occhio umano risolve fino a circa $1'$ (un primo d'arco) al centro del campo visivo (fovea), $\sim1^{\circ}$ in periferia. Un visore VR tipico ha un campo visivo (FOV) di $\sim110^{\circ}$ risolto su $\sim2000$ pixel, dando una risoluzione angolare di $\sim0.05^{\circ}=3'$ -- \textbf{circa 3 volte inferiore} a quella dell'occhio umano.
\end{teorema}

\begin{definizione}[Visione Stereoscopica]
Il senso di profondit\`{a} deriva principalmente dai diversi punti di vista dei due occhi (\textbf{stereoacuity}), nella regione di sovrapposizione dei due coni visivi; \`{e} rafforzato da segnali fisiologici secondari (\textbf{accomodazione}, \textbf{vergenza}) ed \`{e} efficace fino a $\sim$100-400 m, oltre i quali subentrano segnali psicologici (occlusione, dimensione relativa, gradienti).
\end{definizione}

\begin{nota}[Acutezza Stereoscopica -- Test di Howard-Dolman]
Misurando la minima differenza angolare percepibile tra due oggetti a profondit\`{a} leggermente diverse: soglia \textbf{eccellente} $\sim0.1'$, \textbf{normale} $\sim0.5'$ (80\% della popolazione), \textbf{bassa} $\sim2.3'$ (97.3\% della popolazione). La maggior parte dei visori VR ha una risoluzione angolare \textbf{circa 6 volte inferiore} all'acutezza stereoscopica umana media.
\end{nota}

\begin{esempio}[Come un Visore VR Forma l'Immagine]
Uno schermo piatto viene posto a soli $\sim$5 cm dall'occhio (molto pi\`{u} vicino del punto prossimo naturale, $\sim$30 cm): a questa distanza l'occhio non potrebbe mai mettere a fuoco naturalmente. Si interpone quindi una \textbf{lente convergente} (focale $f\sim2.5$ cm) tra schermo e occhio, che forma un'\textbf{immagine virtuale ingrandita} a distanza molto maggiore ($\sim1.2$ m), facilmente osservabile dall'occhio rilassato.
\end{esempio}

\begin{center}
\begin{tikzpicture}[font=\scriptsize, node distance=1.3cm]
\node[draw, circle, minimum size=0.7cm] (occhioL) at (0,0.6) {L};
\node[draw, circle, minimum size=0.7cm] (occhioR) at (0,-0.6) {R};
\node[draw, rounded corners, minimum height=1.6cm, minimum width=0.3cm, fill=blue!10] (lente) at (3,0) {};
\node[below, align=center] at (3,-0.85) {lente\\($f\sim2.5$cm)};
\node[draw, rounded corners, minimum height=1.6cm, minimum width=0.15cm, fill=gray!30] (schermo) at (2.3,0) {};
\node[above, align=center] at (2.3,1.0) {schermo\\($\sim$5cm)};
\draw[-Stealth, dashed, red] (occhioL) -- (7,2.2) node[right]{\textbf{Vergenza} (verso oggetto virtuale)};
\draw[-Stealth, dashed, orange] (occhioL) -- (3.1,0.9) node[right]{\textbf{Accomodazione} (verso piano focale reale)};
\node[align=center, red] at (7,2.7) {oggetto virtuale ($\sim$1.2 m)};
\end{tikzpicture}
\captionof{figure}{Conflitto vergenza-accomodazione: il software regola la parallasse delle due immagini renderizzate per simulare la vergenza corretta verso l'oggetto virtuale (lontano), ma l'accomodazione fisica dell'occhio resta ancorata al piano focale reale della lente ($\sim$1.2 m, fisso) -- una discordanza tra i due segnali di profondit\`{a} che pu\`{o} causare affaticamento visivo.}
\end{center}

\begin{nota}[Il Conflitto Vergenza-Accomodazione]
Nella visione naturale, vergenza (l'angolo con cui i due occhi convergono su un oggetto) e accomodazione (la messa a fuoco del cristallino) puntano sempre alla \textbf{stessa distanza}. In un visore VR con ottica a lente singola fissa, la \textbf{vergenza pu\`{o} essere simulata via software} (conoscendo la distanza interpupillare, IPD, e regolando la parallasse tra le due immagini stereo), ma l'\textbf{accomodazione resta fissa} alla distanza ottica reale della lente ($\sim$1.2 m). Questa discordanza tra i due segnali -- il cosiddetto \textbf{conflitto vergenza-accomodazione} -- \`{e} uno dei limiti fondamentali dell'ottica dei visori VR attuali a lente singola.
\end{nota}

\subsection{Cenni di Ottica Ondulatoria}

\begin{definizione}[La Luce come Onda Elettromagnetica]
La luce trasporta energia in ``pacchetti'' (fotoni): $E_\gamma = hf = hc/\lambda$, con $h=6.63\times10^{-34}$ J s (costante di Planck). Lo spettro \textbf{visibile} \`{e} circa $380$-$750$ nm; oltre si trovano l'infrarosso (IR) e l'ultravioletto (UV). I sensori CCD sono spesso sensibili anche oltre il visibile (fino a $\sim$1000-1100 nm), richiedendo un filtro IR.
\end{definizione}

\begin{nota}[Rifrazione: cambia $\lambda$, non $f$]
Passando da un mezzo a un altro, la \textbf{frequenza} (fissata dalla sorgente) resta invariata; \`{e} la \textbf{lunghezza d'onda} a cambiare, in proporzione a $v$ (e quindi a $1/n$): $\lambda_n = \lambda/n$.
\end{nota}

\begin{esempio}[I Miraggi]
In una giornata calda, l'aria vicino all'asfalto \`{e} pi\`{u} calda (quindi meno densa, $n$ minore, $v$ maggiore) rispetto agli strati superiori: i raggi luminosi si incurvano gradualmente. Il cervello, assumendo (come sempre) propagazione rettilinea, prolunga il raggio percepito in linea retta, facendo apparire un'immagine (es. il cielo, interpretato come una pozza d'acqua) come se provenisse da sotto la strada -- il classico \textbf{miraggio inferiore}.
\end{esempio}

%=============================================================================
\section{Meccanica dei Fluidi}
%=============================================================================

\subsection{Statica dei Fluidi}

\begin{definizione}[Fluido]
Un fluido non ha forma propria (assume quella del recipiente) e \textbf{non sostiene sforzi di taglio}: i \textbf{gas} non hanno n\'{e} forma n\'{e} volume propri (facilmente comprimibili); i \textbf{liquidi} hanno volume proprio e sono difficilmente comprimibili ($B\sim10^9$ N/m$^2$, cfr. \S8).
\end{definizione}

\begin{definizione}[Densit\`{a} e Pressione]
\[
\rho = \frac{m}{V} \qquad\qquad p = \frac{F_\perp}{S}\ \text{[Pa]}
\]
La pressione \`{e} uno \textbf{scalare} (non ha caratteristiche direzionali): su un elemento di fluido in equilibrio, le forze sono sempre normali alla superficie di separazione (altrimenti gli elementi scorrerebbero tra loro). Unit\`{a} comuni: $1\ \text{atm} = 1.013$ bar $= 1.013\times10^5$ Pa.
\end{definizione}

\begin{teorema}[Legge di Stevino]
Per l'equilibrio idrostatico di una colonna di fluido di altezza $h$ sotto la forza peso:
\[
p = p_0 + \rho g h
\]
La pressione dipende solo dalla profondit\`{a} $h$ (non dalla forma del recipiente).
\end{teorema}

\begin{teorema}[Principio di Pascal]
Un cambiamento di pressione esterna $\Delta p_0$ si trasmette \textbf{invariato} a tutti i punti di un fluido incomprimibile, incluse le pareti del recipiente. Conseguenza: le superfici di separazione tra fluidi (o tra fluido e aria) in equilibrio sono sempre \textbf{orizzontali}, indipendentemente dalla forma/inclinazione del recipiente (\textbf{vasi comunicanti}).
\end{teorema}

\begin{esempio}[Pressa Idraulica]
Due pistoni di superficie $S_1,S_2$ collegati da un fluido: per Pascal, $p_1=p_2$, quindi $F_2=F_1(S_2/S_1)$: piccole forze su un pistone stretto generano forze molto pi\`{u} grandi su un pistone largo. Per conservazione dell'energia (fluido ideale, incomprimibile): $F_1d_1=F_2d_2$ (il pistone stretto si sposta di pi\`{u} di quanto si sposti quello largo).
\end{esempio}

\begin{esempio}[Barometro di Torricelli]
Un tubo di vetro capovolto (vuoto all'interno) immerso in mercurio: la colonna di mercurio si alza fino all'altezza $h$ per cui $\rho g h = p_{atm}$, trasformando una misura di altezza in una misura di pressione.
\end{esempio}

\begin{teorema}[Principio di Archimede]
Un corpo immerso in un fluido riceve una spinta verso l'alto pari al \textbf{peso del fluido spostato}:
\[
F_{Arch} = \rho_{fluido}\,V_{immerso}\,g
\]
Un corpo \textbf{galleggia} se la sua densit\`{a} \`{e} minore di quella del fluido ($\rho_o<\rho_f$); \textbf{affonda} se maggiore; resta in equilibrio a qualunque profondit\`{a} se uguale. Un corpo immerso mostra quindi un \textbf{peso apparente ridotto} (es. un oggetto appeso a un dinamometro pesa meno se immerso in acqua).
\end{teorema}

\subsection{Dinamica dei Fluidi}

\begin{definizione}[Viscosit\`{a}]
Attrito interno tra strati di fluido in moto relativo, con coefficiente $\eta$ (\textbf{viscosit\`{a}}): in un fluido reale, gli strati vicini alle pareti di un condotto si muovono pi\`{u} lentamente di quelli centrali (profilo di velocit\`{a} non uniforme), a differenza di un fluido ideale ($\eta=0$) dove tutti gli strati hanno la stessa velocit\`{a}.
\end{definizione}

\begin{teorema}[Teorema di Bernoulli]
Per un fluido ideale (incomprimibile, non viscoso) in moto, lungo una linea di flusso:
\[
p + \rho gy + \frac12\rho v^2 = \text{costante}
\]
equivalente a una legge di conservazione dell'energia (pressione = termine di forza conservativa; $\rho gy$ = energia potenziale; $\frac12\rho v^2$ = energia cinetica, per unit\`{a} di volume). \textbf{Conseguenza controintuitiva}: \textbf{un fluido a velocit\`{a} maggiore ha pressione minore}, e viceversa (es. tetti scoperchiati dal vento, tende della doccia risucchiate verso l'interno).
\end{teorema}

\begin{esempio}[Portanza Alare]
Il profilo di un'ala devia il flusso d'aria in modo che la velocit\`{a} sopra il profilo (dorso) sia maggiore che sotto (ventre): per Bernoulli, la pressione sul dorso \`{e} minore, generando una forza netta verso l'alto (\textbf{portanza}). Data la dipendenza quadratica da $v$, la portanza cresce rapidamente con la velocit\`{a} dell'aeromobile.
\end{esempio}

\begin{definizione}[Flusso Laminare e Turbolento]
\textbf{Laminare}: il fluido scorre in strati regolari e paralleli, senza mescolamento. \textbf{Turbolento}: il moto \`{e} irregolare, con vortici che trasferiscono energia dalla scala macroscopica a quella microscopica (dissipazione). Il \textbf{numero di Reynolds} caratterizza la transizione:
\[
R_e = \frac{\rho v x}{\eta}
\]
Il flusso diventa pi\`{u} turbolento all'aumentare della velocit\`{a} o della distanza percorsa $x$, o al diminuire della viscosit\`{a}.
\end{definizione}

\end{document}
